% Auto-generated by scripts/06_emit_structure.py
% Source: scans 045–073
% Section id: chapter/2
% Phase 3 deliverable. Footnotes (Phase 4), citations (Phase 5),
% and Wade-Giles → Pinyin (Phase 6) are not yet applied here.

\chapter{The Divination Inscriptions}
\label{ch:2}

% source: scan 045, printed 28
% intro-echo-dropped: The Divination
% intro-echo-dropped: Inscriptions
% intro-echo-dropped: 2.I
\section[Introduction]{\origsecnum{2.1}Introduction to the Divination Record}
\label{sec:chapters-ch02:2.1}

Introduction
After the diviner had cracked the oracle bone, the crack was interpreted, and
an engraver then cut a record of the divination into the scapula or plastron. Since we
depend upon that record for our knowledge of what took place, I want next to consider
the record itself, even though it was made after the operations concerned. The record
may be regarded as consisting of a preface (and postface), charge, crack number, crack
notation, prognostication, and verification, though these elements were not all present in
every case. The description that follows refers primarily to the situation in the reign of
Wu Ting (period I); for the way certain of these features changed with time, \ref{sec:chapters-ch04:4.3} (see see sec. 4.3).

\section[The Preface and Postface]{\origsecnum{2.2}The Preface and Postface}
\label{sec:chapters-ch02:2.2}

The preface (called hsü-tz'u by modern scholars) usually recorded the
cyclical day on which the divination was performed, the name of the diviner, and, rarely,
the place of divination. The place of divination, the month, and, by period V, the number
of the sacrificial cycle might also be recorded in a postface after the charge. Presumably,
when the place was not recorded, the divination took place in the ritual center at Hsiao-[\textasciicircum{}1 \textasciicircum{}1]: Though the crack numbers and crack notations were inscribed, I shall, in the interests of clarity, limit the use of the word "inscription" to any combination of preface (and postface), charge, prognostication, and verification, that is, to those units which, singly or in combination, delimit one particular divination topic. The distinction is useful statistically since it enables us to count the number of inscriptions on an oracle bone without including the crack numbers or crack notations. In period I, the full inscription unit of preface (and postface), charge, prognostication, and verification rarely appeared. The most common inscription was either preface and charge or the charge by itself (Li Ta-liang [1972], pp. 108-119, who gives examples of the various combinations of units found in Ping-pien). Also note that the number of inscriptions does not necessarily equal the number of divinations (cf. n. 26), which depends not upon the written record of the divination topic itself, frequently divided into two charges, but on the number of cracks made per topic. One inscription might refer to only one crack, but it might also refer to many (\ref{sec:chapters-ch02:2.5} (see see sec. 2.5)). 4
% source: scan 046, printed 29
t'un. The preface might be abbreviated in various ways, or might in some cases be
omitted entirely.\footnote[2]{It is not established beyond doubt that either the ta-yi Shang, ``the great settlement Shang,'' of the period V inscriptions (S42.3-4) or the Shang where period V divinations were performed (S279.4-280.1) refers to the Hsiao-t'un site; nor is it certain that Hsiao-t'un was indeed the Shang capital (cf. the Preface). For an introduction to the problems, see Ikeda (1964), 1.13.7. On the location of Yin-hsü (or, with GSR 78a, Yin-ch'ü; but cf. Creel [1938], p. 7, n. 11), see Miyazaki (1970; cf. Keightley (1973), p. 537, nn. 36, 37.}\footnote[3]{Preface abbreviations have been treated exhaustively by Chou Hung-hsiang (1969), pp. 60-67, and Li Ta-liang (1972), pp. 120-126.}\footnote[4]{Not infrequently, the preface, in whole or in part, might appear on the front of the shell and the charge on the back, or vice versa (Li Ta-liang [1972], pp. 70-71, 74-79; Ping-pien 358/359 is a good example). Cf. n. 65 below.}\footnote[5]{Ping-pien 1.12 (period I).}\footnote[6]{The bone graph for this diviner's name is generally transcribed as (e.g., Ping-pien 1, k'aoshih, p. 14), which should be read Ch'üeh. See Takashima (1973), p. 324, n. 7; GSR 1226a. Some scholars, however, read the name K'o (e.g., Jao [1959], English table of contents, p. 2; Serruys [1974], p. 43). As Takashima notes, the assignment of modern Mandarin readings to Shang graphs is largely a matter of convenience; in many cases, especially for names, the romanization must be based upon what is supposed to have been a phonetic element in the character concerned (Barnard [1960], p. 105).}\footnote[7]{Traditionally, following the Shuo-wen definition, cheng pu wen yeh, “Cheng is to question by crack-making,'' the M has been translated as ``to question,'' with the charge posed as an interrogative (e.g., ``Today will it rain? / Today will it not rain?'' [Ping-pien 263.9–10]). But it is not certain that oracular (a simplified picture of a cauldron) should be read as cheng ti rather than ting (According to Karlgren, cheng l'i had the archaic reading of *tieng; ting had the archaic reading of * tieng [GSR 834a,g]; Chou Fakao et al. [1973], 9284, 11944, gives *tieng and *teng respectively). And, more importantly, I believe that on the basis of both etymology and Shang usage, was not interrogative and that the charge should be translated as a prayer, prediction, or statement of intent (e.g., ``Today it will perhaps rain / Today it will not rain''). The evidence for this, initially presented in Keightley [1972], will be fully discussed in Studies.) Takashima \listitem{1973} and Serruys \listitem{1974} also take as an indicative rather than interrogative verb. and translate it as ``test.'' The precise meaning of is not yet clear to me, and I prefer to use ``divine'' as an interim, functional translation; and since we are not sure which, if either, word, cheng or ting, was its primary descendant, I prefer to romanize the graph as *tieng, Karlgren's archaic pronunciation for both when the diacritical marks are omitted.}\footnote[8]{Chin-chang 583 (D) (S201.4; period V..}\footnote[9]{Ch'ien-pien 2.17.3 + 2.17.5 (S202.2; period V). These scapula fragments have been joined with others to form Ho-pien 218 (D) (=Chui-hsin 345). The words ``it was'' may be supplied on the basis of standard Shang dating practice in period V inscriptions on both bronze and bone e.g., Matsumaru [1970], pp. 65-68); the catch of foxes was not just a record of fact but a way of specifying which first month it was. For the place names (and on their pronunciation, cf. n. 6), see CKWT, pp. 3385, 3392, 4125. Tung (1945), pt. 2, ch. 9, pp. 61a-63b, discusses the geography; cf. Ikeda (1964), 1.15.12, 1.15.14. BA-x1 E (綴二十八) 29 U}
% source: scan 047, printed 30
The cyclical dates recorded in the prefaces reveal that divination took place on every
day of the ten-day Shang week. Certain regularities of schedule may be discerned,\footnote[10]{The following regularities may be cited: \listitem{1} Kuei-day divinations were more numerous than divinations made on any other day (Durrant [1972], p. 9). This is partly because routine divinations of the form ``Crack-making on kueiyu, Cheng divined: 'In the (next) ten days there will be no disaster'''' always took place on the kuei-day, immediately preceding the chia-day which started the new ten-day week (sec. 2.3.1, par. 5). For an abbreviated listing of such divinations, see S165.1-168.2. Chia-t'u 80, a fragmentary plastron from period I, contains twenty-three divination charges of this type, ranging from the day kuei-yu of the tenth month of one year to the day kuei-hai of the fifth month of the next year. \listitem{2} Preliminary research has indicated that divinations about certain rituals like the liao (S202.4), chiu (S391.2-395.1), and yi 1 (S250.2-251.2) tended to take place more frequently on kuei-days than on any other (Durrant [1972] \listitem{3} Since the sacrifices in the regular five-ritual cycle were offered to ancestors on their cyclical name-days, divinations about sacrifice usually took place either on the name-day itself or on the day prior to the sacrifice. As the sacrificial schedule became more regular (appendix 5, sec. 2), so did the divinations about it. \listitem{4} Matsumaru \listitem{1963} has found that hunt divinations were made according to no fixed schedule in period I; at a certain point in period II, hunt divinations started to be limited to three days (yi, wu, and hsin) of the ten-day week; in III+IV, the jen-day was added; and in period V, the ting-day was added, making a total of five hunt days in all. (The validity of his conclusions may readily be checked by surveying the approximately 1,680 hunt inscriptions listed at S289.3-298.2.)} 10
but every day was likely to involve a stream of pyromantic burnings and crackings.
Preliminary research indicates no evident periodicity in the daily work schedules of
particular diviners.\footnote[11]{This conclusion is based upon all the divinations made by the period I diviner Ch'üeh that are listed in Jao (1959), pp. 73-240; I am grateful to Inge Dietrich for making this study. More information about work schedules might be discovered by studying the more systematic divinations of period IIb. Tung's hypothetical reconstruction of five scapulas ([1945], pt. 2, ch. 6, pp. 6b-7b), for example, implies that the IIb diviner Lü divined about every night for a two-month period.}11
I can only speculate about why the preface was recorded. Presumably, the reasons
were both religious and historical. The date of the divination recorded the time when
the spiritual forces of the universe, revealed by the cracks, had been identified. Further,
the cyclical date established that sacrifices to the ancestors-who were named by the
cyclical stems Father Chia, Father Yi, and so on—had been offered on the appropriate
name-day; the prefaces thus documented the orderly execution of the ritual schedule
upon which the goodwill of the ancestors depended. Similarly, the carving of the diviner's
name would have served to identify the officer responsible for overseeing the divination
ritual, further documenting the full nature of each particular contact that was established
between man and the spirits.\footnote[12]{The metaphysics and function of Shang divination will be discussed in Studies; for a preliminary exploration, see Keightley (1973a) and (1975c). It is conceivable that the record of the diviner's name would have permitted the king to evaluate a diviner's performance, but the small number of verifications (sec. 2.8)—those records that would have been required to evaluate the success or failure of each prognostication (sec. 2.7) -make this unlikely. The fact that the verifications we do have generally prove the forecast to have been correct again suggests that the verifications were not carved to document the success rate of particular diviners; divination failures were simply not recorded on the bone or shell. On this point see n. 83.} 12
C
% source: scan 048, printed 31

\subsection{\origsecnum{2.2.1}The Diviners}
\label{sec:chapters-ch02:2.2.1}

The Diviners
The names of approximately one hundred and twenty diviners have been
identified on the oracle bones,\footnote[13]{Ch'en Meng-chia (1956), p. 202, gives a total of 120 diviners. Shima (1958), p. 34, lists III diviners; S557-576 lists 104 diviners: Jao (1959), p. 1202, gives a total of 117 diviners, with 20 more names needing further study. The discrepancies, which generally involve unique cases and disputed readings, are not significant. For specific examples of disagreement, see Li Hsueh-ch'in (1957), pp. 125-126, who would delete 21 of the names proposed by Ch'en Meng-chia and add 5 new ones; Yen Yi-p'ing (1961), p. 521, who would delete I and add 4; see too the note to table 6.} 13 but the bulk of the divinations for which a diviner's
name is recorded was performed by five period I diviners (Ch'üeh, Pin, Cheng,\footnote[14]{Ch'ü Wan-li (Chia-pien 138, k'ao shih, p. 22) reads this diviner's name, as Yi; for rea, sons which support the more generally accepted reading of Cheng, see Serruys (1974), p. 86, n. 3.}14
Hsüan E, and Chung [? S557-563]), six period II diviners (Ta, Lü, Hsing,
Chi, Yin, Ch'u [S566-570]), and one diviner from period III (Ho [S572]).
Diviners' names are a valuable dating criterion (sec. 4.3.1.2). The fact that their names
were generally recorded in periods I and II but are rarely recorded in IV and V suggests
that their role became less important. By period V, the king-who served as diviner in
all five periods was virtually the only diviner and prognosticator recorded.\footnote[15]{See table 7, nos. 1.2.2, 1.5.1, 1.5.2, 3.2, etc. The way in which diviners' names may be used to indicate the evolving status of the diviners, and thus of divination itself, in Shang society will be considered in Studies.} 15

\subsubsection[Diviner Groups]{\origsecnum{2.2.1.1}Diviner Groups}
\label{sec:chapters-ch02:2.2.1.1}

The Shang diviners may be divided into various groups. They may be separated
primarily by the periods in which they were active, so that one may speak of the period
I diviners by contrast with those, for example, of period IIb. But within the periods
themselves, subgroups of diviners may be discerned, primarily by the appearance of
their names on the same bone, by the similar content of their divinations, and by common
pit provenance.\footnote[16]{The hypothesis that two or more diviners divined a single charge jointly has been effectively refuted by Li Yen (1972).} 16 In period I, for example, most divinations for which a diviner's name
was recorded were conducted by approximately twenty-five ``court diviners.'' These
diviners called the Pin group by Ch'en Meng-chia, after the diviner whose divinations are the most numerous and typical of the group-dealt with matters involving the
state as a whole: weather, agriculture, warfare, sacrifices to the nature powers, sickness
of the ruling elite, and so on. The ancestral sacrifices with which they were concerned
were generally limited to those addressed to dynastic ancestors (and their consorts) who
had ascended the throne; it is for this reason that this group may be regarded as court
diviners, concerned with sacrifices offered to members of the ta-tsung, the great
ancestors who had actually ruled the Shang state.\footnote[17]{Ch’en Meng-chia (1956), p. 148; cf. Ch’iu (1972), p. 43. I derive the term ``court diviners'' from the discussion of Kaizuka (1946), p. 236. On the ta-tsung, see \ref{ch:4} (see ch. 4), n. 18. 31 || ||} 17
Other groups of diviners who were practicing, either in period I or period IV (their
% source: scan 049, printed 32
date is in dispute), are known to modern scholars as the Royal Family group and the
Tzu group. The rituals about which they divined included sacrifices addressed to
ancestors and their consorts who had ruled; but modern scholars associate these diviners
with the royal family because their divinations paid particular attention to sacrifices
offered those classificatory ``fathers,” “mothers,'' ``brothers,'' and ``sons'' of the royal
family who had not ruled.\footnote[18]{The periodization of diviners of the Royal Family and Tzu groups see table 6 for their names) has been a matter of intense controversy. Tung Tso-pin and his supporters (e.g., Li Hsüehch'in [1957]; Shima [1958]; Yen Yi-p'ing [1961]; Hsü Chin-hsiung [1973] and [1974]; cf. Chang Ping-ch'üan [Ping-pien 611. k'ao shih, p. 95]) would place them in period IVb. Scholarly opinion, however, has sided with Tung's critics, notably Kaizuka and Itō \listitem{1953} and Ch'en Meng-chia (1956), pp. 145-155, 158-167 (who divides the disputed inscriptions into the three diviner groups of Tui, Tzu f, and Wu F), placing the disputed inscriptions in period I or early II; e.g., Jao (1959), pp. 657-658, 678, 687, 747, and the summary at pp. 1204-1210; cf. Jao (1961a), p. 95; Matsumaru (1963); Ikeda (1964), preface, p. 2, and passim; Chang Tsung-tung (1970), pp. 23, 27, and passim; Shirakawa (1972), p. 125; Maekawa (1974). Hu Hou-hsüan provisionally placed such inscriptions (Ching-chin 2908-3068; 3115-3160; Hsü-ts'un, 1.1433-1.1465; 2.579-2.598) prior to period (Ching-chin, preface, p. 1b; Hsü-ts'un, preface, p. 1). Ch'ü Wan-li (Chia-pien k'ao-shih, preface, pp. 3-4, 9) places Tung's period IV inscriptions in period I, though he admits there may be some errors. In this book, I use the abbreviation RFG as a generic term for any or all of these various subgroups. The issues are too technical to be resolved here; a detailed discussion of the evidence and the reasons why I would place the disputed inscriptions in period I will be presented in Studies. For a recent archaeological find which may support an early date for the RFG diviners, see \ref{ch:4} (see ch. 4), n. 179. Note that if the RFG inscriptions are dated early, Tung's hypothesis ([1945], pt. 1, \ref{ch:1} (see ch. 1), p. 2b; pt. 2, \ref{ch:5} (see ch. 5), p. 2b; [1964], pp. 89-106) that a so-called Old School and New School (cf.\ref{sec:chapters-ch04:4.2} (see  sec. 4.2)) alternately held power and shaped religious and institutional practices is without foundation. This is not to deny that a New School appears to have taken over in period IIb (see sec. 4.3.1.12); it is the hypothesis of the Old School's return in IVb that would be invalidated; the disputed inscriptions which show Old School characteristics would do so because they were contemporary with the Old School of period I, not because they were a period IVb revival of it.} 18 Other groupings of diviners may have existed in other periods,
but more research is needed to document, and provide a rationale for, their existence.\footnote[19]{Itō (1962), pp. 19-21, 35, n. 46, has suggested, for example, that the Ta group of diviners in period II, headed by the diviner Tax, tended, like the Old School diviners (see n. 18 and\ref{sec:chapters-ch04:4.2} (see  sec. 4.2)) of period I, to divine about the worship of local nature powers and political activities of interest to the surrounding statelets. By contrast, the Yin group of period II diviners, as represented by the diviner Hsing f, divined about the worship of the former kings of Shang and about political matters that directly affected king Tsu Chia, the IIb ruler. Itō traces this difference in divination focus to the origins of the two groups. He suggests that the diviner Hsing was a member of a group that had had close relations with the Shang royal house from period I on; hence the more central nature of its concerns. The diviner Ta, by contrast, came from a statelet or localized lineage that had been conquered by the Shang in period I; hence the more provincial nature of this group's divinations (cf. Itō [1961], pp. 283-284). Matsumaru (1963), pp. 62-66, 72, n. 6, also discerns the existence of two groups of diviners in the period IIb hunt divinations. His A group divined according to no particular schedule and with other miscellaneous topics engraved on the bone; his B group divined according to a schedule and reserved the whole bone for hunt inscriptions alone. He regards the B group inscriptions as prototypes for the orderly inscriptions typical of period V (see sec. 4.3.1.5; appendix 5, sec. 2). These hypotheses are possible, but most of the distinctions cannot be drawn rigidly. Until further evidence is available, it is best to think in terms of tendencies rather than absolute divisions.} 19
% source: scan 050, printed 33

\section[The Charge]{\origsecnum{2.3}The Charge}
\label{sec:chapters-ch02:2.3}

The term ``charge'' (ming-tz'u) is used by modern scholars to refer to the
topic of the divination inscription. There is no certain evidence that ming was a Shang
term. 20 That it was used in Chou plastromancy, however, suggests that the modern
term may have been used in Shang times; the Chou usage supports the view that the
individual divination sentences were not questions but were wishes, tentative forecasts,
or statements of intent with which the bone or shell was being charged. 21 The charge was
usually, though by no means invariably, carved on the front of the scapula or plastron.\footnote[22]{Cf. n. 65.} 22

\subsection[Content]{\origsecnum{2.3.1}Content}
\label{sec:chapters-ch02:2.3.1}

Few significant aspects of Shang life were undivined (on this point,\ref{sec:chapters-ch05:5.2} (see  cf. sec. 5.2)).
The following list of topics, by no means exhaustive and frequently overlapping, gives
the major areas of concern in period I:\footnote[23]{For other listings, see Tung (1965), p. 116; Chang Ping-ch'üan (1967), p. 857; cf. n. 26 below. Sample inscriptions may be consulted as follows: \listitem{1} sacrifices (S91.3-92.3); \listitem{2} military campaigns (S329.4-331.1); \listitem{3} hunting expeditions (S289.3298.2); \listitem{4} excursions (\$67.3-68.3; 79.1); \listitem{5} ten-day week (S165.1-166.3); \listitem{6} night or day (S161.4-162.5; 105.1-2); \listitem{7} weather (S169.3170.2); \listitem{8} agriculture (S193.1; 195.1-196.4; 197.2-198.2); \listitem{9} sickness (S447.2–448.2); \listitem{10} childbirth (S281.4-282.3); \listitem{11} distress or trouble (S303.1-305.4); \listitem{12} dreams (S450.3-451.4); \listitem{13} settlement building (S43.2-3); \listitem{14} orders (S46.3-4); \listitem{15} tribute payments (see \ref{ch:1} (see ch. 1), n. 41); \listitem{16} divine assistance or approval (S47.2-3; 463.2-4); \listitem{17} requests (S192.1-4).} 23
I. Sacrifices. The divinations generally sought to establish in whole or in part if
the offering of a particular ritual, with a particular kind or number of sacrificial victims,
to a particular ancestor or group of ancestors on a particular day would be without fault
or without blame and thus auspicious and efficacious.
2. Military campaigns. The major thrust of these divinations was to determine if
Ti, the high god,\footnote[24]{Ti had dominion over rain, wind, other atmospheric phenomena, harvests, the fate of urban settlements, warfare, sickness, and the king's person. He shared some of these jurisdictions with the ancestors and nature powers, but he was the supreme religious deity.} 24 would assist the Shang, if the king should ally himself with another
leader or statelet, if he should issue orders to a particular officer, if he should mobilize
so many thousands of men, if he should send them first to one place and then to another,
and the like.
20. The crack notation may possibly be
read as pu tsai ming “not again charge (the
oracle bone).'' See \ref{ch:4} (see ch. 4), n. 134.
21. The phrase ling-kuei , or ming-kuei,
``to charge the turtle,'' appears in Shang-shu, ``Chin
t'eng,'' and in Tso-chuan, Chou-li, and Li-chi. For
the noninterrogative nature of the Shang charges,
see n. 7. As indicated in the preamble to this book,
it is possible that the charge was spoken at the
time of divination; see, for example, the discussion
of the preface formula wang yüeh *tieng,
``the king spoke and divined,'' at Jimbun 1460,
shakubun, p. 405, n. 1; cf. Eisenberger (1938), p.
109. It is also possible, at least for some divinations,
that the charge carved on bone was only an
abbreviation of a more complete charge proposed
orally for each crack (cf.\ref{sec:chapters-ch03:3.7.4.2} (see  sec. 3.7.4.2)); this will be
discussed in Studies.
% source: scan 051, printed 34
3. Hunting expeditions. These divinations dealt with such matters as the site of
the hunt, its possible success, the possibility of bad weather, and the desirability of avoiding
misfortune en route. They frequently end with the incantation ``there will be no disaster.''
4. Excursions. These divinations sought to establish if the king should travel to a
certain place, go to inspect a certain activity, and so forth and were again concerned with
the desirability of avoiding misfortune en route. The standard incantation is ``going and
coming there will be no disaster.''
5. The ten-day week (hsün ). These divinations, which were always made on the
kuei-day, the last day of the week, sought to establish by means of routine, ritual incantation
that ``in the (next) ten days there will be no disaster.”
6. The night or the day. Like the ``ten-day'' divinations, these charges sought to
establish that the whole day or the coming night would be without disaster. Once again,
I regard these divinations as routine incantations to ward off misfortune and danger. Some
appear to have been made when the king was traveling, hunting, or making war and was
without the defenses, spiritual as well as military, of his capital.
7. The weather. In addition to divinations about the weather to be encountered
while hunting or traveling, the Shang sought to determine if it would rain, the timing of
the rain and its quantity, if the rain would stop, if the sun would shine, if there would
be winds, or snow or other forms of precipitation. The concern was to discover how the
weather might affect agriculture, ritual activity, or the king's various expeditions.
8. Agriculture. In addition to divinations about rainfall, the Shang also divined
about working the fields, about planting grain, about ordering officers and mobilizing
labor gangs to reap and sow, about which spirit might be cursing the crops, and, above
all, about whether Ti would confer harvest.
9. Sickness. These divinations generally specified the nature of the ailment―
toothache, sick eyes, pain in the foot-and sought to determine which ancestor was cursing
the king in this way and what sacrifice might effect a cure.
IO. Childbirth. The charges sought to determine whether or not the childbearing
of a particular royal consort would be “good,'' that is, whether the offspring would be
male, and if the birth would occur without misfortune.
II.
Distress or trouble. These divinations sought to establish whether or not the
king, or a member of the royal family, or one of the king's allies would suffer distress or
trouble of some unspecified sort.
12.
Dreams. The charges sought to determine the potentially ominous significance
of a dream, usually the king's, that had already taken place and to discover which ancestor
or power had caused the dream.
13. Settlement building. The charge indicated that an urban settlement was about
to be built and sought to determine whether or not Ti approved the decision.
14. Orders. The charge generally stated that a particular officer was being ordered
to perform some task; an auspicious divination presumably validated the wisdom of the
choice.
% source: scan 052, printed 35
15. Tribute payments. 25 The charges sought to determine if a certain chief or
statelet would be sending in horses, dogs, bovids, turtle shells, or Chiang captives to the
Shang and, on occasion, in what quantity.
16.
Divine assistance or approval. Divinations about a wide variety of topics ended
with the wish, frequent in the case of military campaigns against enemy statelets, ``We will
receive assistance'' or ``There will be approval.'' The approval sought was that of Ti and
the ancestors.
17. Requests addressed to ancestral or nature powers. These divinations sought to
determine if ritual requests or prayers for progeny, harvest, rain, assistance, and the like
would be auspicious, that is, successful. The charges were frequently concerned with
determining to which ancestor the ritual ought to be addressed.
It should be possible, using Shima's Inkyo bokuji sōrui (sec. 3.3.2), to calculate the
relative frequency of these various divination topics, but such a comprehensive study has
not yet been made. 26 There is no doubt that in periods I, II, and V divinations about
sacrifice were far more numerous than those about any other topic. 27 In period I, divinations about enemy attack, weather, the ten-day week, the issuing of orders, and so forth
were also present in significant numbers. By period V, the scope of Shang divination had
dwindled remarkably and was concerned primarily with divinations about sacrifice, the
ten-day week, and the hunt (sec. 4.3.1.12).
Temporal Range
Many divinations about sacrifice appear to have taken place either immediately
before or simultaneously with the offering of the sacrifice itself. 28 A few divinations, such
as those which sought to identify the power responsible for a dream, sickness, or crop
failure, referred to the immediate past and present. The bulk of the divinations about other
25. But cf. \ref{ch:1} (see ch. 1). n. 39.
26. Kaizuka (1947), p. 468, revising Lo Chen-
yü's original statistics, gives a table based on 1,203
inscriptions; 538 deal with sacrifices, 196 with
hunting and fishing, 34 with harvests, etc; see too,
Li Chi (1977), pp. 28-29. Such counts, however,
need to separate inscriptions by periods and must
be based on a far larger sample. Further, all such
counts refer only to the number of inscriptions,
not the number of divinations (see n. 1), so great
care is needed in assessing the relative significance
the Shang attached to certain topics (cf. Chang
Ping-ch'üan [1956], pp. 247, 256-257).
27. I am uncertain about the situation in periods III+IV, for which a large number of hunt
divinations (S290.1-293.4), but not many sacrifice
divinations, has been found. We cannot be sure
whether this represents a real shift in Shang divination concerns or whether it is due to the accidents of retrieval (cf.\ref{sec:chapters-ch05:5.3.1} (see  sec. 5.3.1)). It seems unlikely,
however, that the Shang would have remained
sufficiently ``superstitious'' to make numerous hunt
divinations without remaining equally ``superstitious'' about the need to divine about sacrifices
to the powers; it was the powers, after all, who
validated the entire divination process (the theology of Shang divination will be considered in
Studies).
28. E.g., Chui-hsin 304 (fig. 7), translated in
appendix 5, sec. 2. For a preliminary study of the
virtual simultaneity in many cases of divination
and sacrifice, see Keightley (1972), pp. 35-37:
The assumption is that the divination determined
the acceptability of the sacrifice as it was, or was
about to be, offered (cf. appendix 5, n. 6). The
issue will be discussed further in Studies.
% source: scan 053, printed 36
topics, however, was concerned with the immediate future-the night, the morrow,
or the coming ten-day week. 29
More rarely, the diviner would look further ahead. Period I divinations about sacrifice
might be made, for example, eight, nine, eleven, fourteen, or twenty days in advance, 30
reflecting a time when the sacrificial schedule was not yet firmly established and it was
presumably still necessary to determine an auspicious day before the offering. In period
IIb, and especially in period V, when the sacrificial schedule had become regular (appendix
5, sec. 2), such long-range divinations were no longer needed. A similar loss of interest in
the long-range view may, in fact, be discerned in all areas of Shang divination. With the
advent of the New School in Period IIb, the diviners appear to have concentrated upon
the immediate spiritual climate of the day, the morrow, and the ten-day week. They were
no longer concerned with long-range forecasts or prognostications. 31 These trends may
be related to a general reduction in the scope of Shang pyromancy (sec. 4.3.1.12).

\section[Crack Numbers]{\origsecnum{2.4}Crack Numbers}
\label{sec:chapters-ch02:2.4}

After the Shang diviners had formed cracks on the front of the bone or shell,
they frequently, though by no mean always, numbered the cracks to indicate their sequence
in the set (on sets, \ref{sec:chapters-ch02:2.5} (see see sec. 2.5)). 32 In period I, the crack number (known by the modern
terms pu-chao chi-shu-tz'uor hsü-shu ) 33 was a number from 1 to 10.34
It was incised into the front (never the back) of the scapula or plastron before the charge
itself was carved. 35 It was generally placed to the right or left of the upper end of the vertical
crack, on the same side, right or left, as the transverse crack (fig. 8); occasionally, it was
29. Period I divinations enquiring whether it
would rain ``from today until a future kan-chih
day” (S261.1-3) span, on the average, a period
of four days. For a typical divination about sacrifice on the next day, see Chien-shou 5.4, translated
in appendix 5, sec. 2.
30. Ping-pien 207.4; 209.3; 207.3; 207.1; 197.3-
4. On the increasing certainty of the ritual schedule in period IIb, see appendix 5, sec. 2.
31. The rainfall divinations of the type referred
to in n. 29 were no longer performed after period
II. Prognostications which, in period I, might be
both long-range and specific as to the day (sec.
2.7.1) ceased to be either in later periods.
32. Chang Ping-ch'üan (1956), p. 237, notes
that numberless cracks are frequent on turtle shells,
whether inscribed or not; that many cracks with
crack numbers lack inscriptions, but that in some
cases the inscriptions may be found on the other
side of the shell or even on another shell (cf.
Ping-pien 24, k'ao-shih, pp. 51-52; Yen Yi-p'ing
[1973], p. 7a). For examples of numberless cracks
-which can sometimes be located in rubbings by
making a tracing of the hollow positions and
superimposing the tracing paper (reversed) on the
front of the shell to see where cracks could be
expected to have formed-\ref{sec:chapters-ch03:3.7.2} (see see sec. 3.7.2); cf.
Ping-pien 62, k'ao-shih, p. 96. The reason why
certain cracks were not numbered is unclear (cf.
n. 63).
33. Tung (1945), pt. 2, \ref{ch:3} (see ch. 3), p. 26a; Chang
Ping-ch'üan (1954), p. 226.
34. Yen Yi-p'ing (1961), p. 486, has suggested
that crack number 11 appears twice on Yi-pien
(fig. 5); if this is so-and study of the actual@@
plastron would be needed to determine if these
marks are true crack numbers-it is a unique
case.
35. I see no way to tell whether, as Tung claims
([1949a], p. 260), the crack numbers were written
after the cracks had been read, rather than before.
Nor can we tell whether the crack numbers were
cut after each single cracking, as Chang Ping-ch'üan supposes ([1956], p. 230), or whether they
were cut in sequence after all the cracks in a set
were made.
% source: scan 054, printed 37
placed at the end of the transverse crack.\footnote[36]{Chang Ping-ch'üan (1954), pp. 226, 227; (1956), p. 230; (1967), p. 861. On a few RFG plastrons the crack number is placed to the left or right of the bottom end of the vertical crack, as in Yi-pien 1318 (Chang [1954], p. 226; [1956], p. 230).}36 Crack numbers, as we have seen (sec. 1.6.3),
enable us to determine the order in which the hollows were burned; they assist the reconstruction of scapulas and plastrons from bone and shell fragments (sec. 5.7); and they are
important for the study of sets (sec.\footnote[37]{Chang Ping-ch'üan (1956), pp. 240-241. The location of the crack number indicates the direction of the transverse crack, which in turn may indicate whether the fragment comes from the right or left side of a turtle shell (on the patterns made by the cracks, see secs. 1.5.1; 1.6.2; 1.6.3).} 2.5).37
When the grooves of an engraved graph actually straddled a crack (a phenomenon
given the modern name fan-chao\footnote[38]{Chang Ping-ch'üan (1954), p. 230. For examples of fan-chao, see nn. 39, 87.} ), 38 the crack numbers were commonly erased and
sometimes placed elsewhere. It is this fact which tells us that the crack numbers were
engraved before the charge was carved.\footnote[39]{Chang Ping-ch'üan (1954), pp. 226, 230; (1956), pp. 230-231, 236, 237; cf. Ping-pien 49, k'ao-shih, p. 83. Chang notes that these erasures are common but not always easy to discern in rubbings. He cites as examples Yi-pien 1354; 1500; 1905.}39 Why most straddled cracks were denumbered
in this way is not clear; the crack numbers may have been erased to avoid any confusion
with the content of the inscription.\footnote[40]{Chang Ping-ch'üan (1954), p. 230. This proposal is of interest, for it suggests that somebody, either person or spirit, who had not been immediately involved in the divination itself-or, if he had been, had forgotten the details of the charge was expected to read the engraved inscriptions.} 40
The routine use of crack numbers in period I is documented by their appearance on
the front of every reconstituted Ping-pien plastron. We know it was important to record
the crack numbers accurately because wrong crack numbers were frequently erased and
corrected;\footnote[41]{E.g., Ping-pien 47.15; 49.11; 93·4; 134·5.}41 on the other hand, wrong or omitted crack numbers were on occasion allowed
to pass,
an understandable lapse in view of the many crack numbers carved. The attention
the Shang paid to crack numbers suggests that a particular crack's place in the sequence of
cracks which formed a set affected the way the cracks were prognosticated.\footnote[42]{E.g., Ping-pien 39.21; 43.5; 43.15; 81.5; Yipien 3398.}43\footnote[43]{Cf.\ref{sec:chapters-ch03:3.7.4.2} (see  sec. 3.7.4.2). How the cracks were read will be discussed in Studies.}

\section[Divination Sets]{\origsecnum{2.5}Divination Sets}
\label{sec:chapters-ch02:2.5}

The response to divination charges was not, in many cases, obtained from only
one crack in one hollow, but from a series of cracks in a series of hollows. Modern scholars
refer to such a series as a ``set'' (ch'eng-tao ). A set may be defined as the series of divinations represented by one or more sequences of numbered cracks, all made on the same
day, with reference to the same topic.\footnote[44]{Chang Ping-ch'üan (1960); (1967), pp., 868-870, provides a good introduction to sets, as does his k'ao-shih to the sets or partial sets in Ping-pien (examples of which are conveniently assembled by Li Ta-liang [1972], pp. 100-107). Two or more sequences of numbered cracks were involved when a single divination topic was divided into two or more charges. For useful quali32 37 ] | | | D}44 Divination topics were frequently divided into a
% source: scan 055, printed 38
pair of charges (modern term tui-chen) in the positive and negative mode; for example,
``We will receive millet harvest/We will not perhaps receive millet harvest.”\footnote[45]{Ping-pien 8 (fig. 8). The significance of such positive and negative charges will be considered in Studies; for a preliminary assessment, see Keightley (1973a), pp. 13-19, 37-38, 45-49; (1975), pp. 17-24; see too n. 124. As Li Ta-liang (1972), pp. 104-107, notes, not all sets on shell were evenly divided; some might have positive charges on both right and left, others might have two positive charges on one side to only one negative charge on the other (or vice versa). Similar variation may be found on scapulas on which above-below opposition of pairs was more usual than the right-left opposition generally found on plastrons (Chang Ping-ch'uan [1967], p. 868); for numerous examples of positive-negative and other types of pairs on scapulas, see Chou Hunghsiang (1969), pp. 37-52; cf. Tung (1945), pt. 2, \ref{ch:5} (see ch. 5), p. 22.} 45 In such
cases, five cracks (numbered 1 to 5) might be made for the positive charge and five more
cracks (also numbered 1 to 5) made for the negative charge, producing a set of ten cracks
in all for this one topic (fig. 8). Generally speaking, the Shang diviners preferred making
cracks in multiples of five, but in period I there appears to have been no standard number
of cracks per set. (For the situation in later periods, see sec. 4.3.1.10.) I would propose the
following classification for sets on period I plastrons: \listitem{1} regular sets, with five cracks on
the positive side and five cracks on the negative, making ten cracks in all;\footnote[46]{E.g., Ping-pien 5; 8 (fig. 8); 12-21 (see sec. (3.7);34-38;174.} 46 \listitem{2} extra-large
sets, with ten cracks per charge, making, in the case of paired charges, twenty cracks per
topic;\footnote[47]{E.g., Ping-pien 49.1-4; 67.7-8; 108.5-6; 145.1-2; 147.1-4; 169.5-6; 172.2-3; 172.6-7; 199.11-12; 199.17-18. (The presence of crack number 11 [n. 34] would raise the possibility of a twenty-two-crack set.) It might be thought that such extra-large sets were the result of a failure of the first ten cracks to yield a reading (the twentytwo crack set, for example, lacks any crack notations), but that is clearly not the case. Ping-pien 108.5 has a shang-chi crack notation by crack 1, yet nine more cracks were made for that charge. Ping-pien 126.10 has a shang-chi by crack 1, yet six more cracks were made for that charge. This suggests the opposite alternative, that large sets were the result of a desire to pursue a ``lucky streak,'' in the hopes of more than one auspicious reading per set. Thus, Ping-pien 145.1-2 has three shang-chi and two hsiao-chi crack notations for the twenty cracks made; the fact that four of these five auspicious readings had presumably been obtained by the time the first ten cracks had been formed (cf. Ping-pien 172.6-7) may have encouraged the diviner to crack ten more in the hopes of a ``jackpot'' set. According to such a hypothesis, we would expect a high proportion of the extra-large, ten-plus crack sets to have auspicious crack notations early in the series; preliminary research, however, has not supported such a conclusion.} 47 \listitem{3} medium-large sets, which appear to have had from five or six to nine or ten
cracks per side, making, in the case of paired charges, a total of eleven to nineteen cracks
per topic;\footnote[48]{E.g., Ping-pien 41.16–17; 49.10-13;67.5-6; 67.9-10; 108.3-4; 126.10-11; 134.1-2; 145.5-8; 147.7-8; 149.11-12; 151.1-4; 157.1-2; 169.1516; 172.10-11; 189.1-2; 199.15-16; 264.7-8. No set with over five cracks per side ever continues on a second plastron (see n. 52). Why some sets stopped between five and ten cracks per side is not clear. In some cases, it seems that the diviner, disgusted with a set which had yielded no response at all after thirteen or more cracks, simply gave up without completing the set and started afresh, e.g., Ping-pien 151.1-2 (nine and seven cracks); 157.1-2 (eight and nine cracks); 189.1-2 (seven and six cracks); 264.7-8 (eight and eight cracks).} 48 \listitem{4} extra-small sets, which contain less than five cracks to a charge or less than
ten cracks to a pair of charges.\footnote[49]{Although we cannot be sure that the set} 49
fications of Chang's definition of a “set,'' see Li
Ta-liang (1972), pp. 100-103. It is highly probable
that further sequences of numbered cracks were
made for some topics, but with no new inscription
recorded (e.g., Ping-pien 49.3-6; 49.8-9; 49.12-
13; 61.7-11). These inscriptionless sets, which
suggest that certain topics might have been cracked
forty times or more, are not included in the discussion that follows since their topic cannot be
identified with certainty.
% source: scan 056, printed 39
The cracks of a set might all appear on one bone or shell or on five separate bones or
shells.\footnote[50]{All on one shell: Ping-pien 8 (fig. 8); 41.1617; 49.10-13; on five separate shells: Ping-pien 12-21 (see sec. 3.71; 34-38 (note that Ping-pien 385 is a more fully restored version of Ping-pien 38); 71 and 73; Ping-pien 80 and Yi-pien 2235; Ping-pien 207 and 209. Li Ta-liang (1972), pp. 166-169, describes the multiple-plastron sets in Ping-pien. As Chang Ping-ch'üan (1956), p. 253, notes, the five plastrons forming a set were all chosen to be about the same size; they were also probably of the same species (i.e., Ping-pien 34; 36; 37: 38 [=385] are identified as Os). For the possibility that scapulas or plastrons may also have been divined in groups of three, see Kuo Mo-jo (1972), pp. 2, 3; Hu Hou-hsüan (1972a), pp. 5-6; but cf. the response by Ch'iu (1972), PP. 43-44.}50 If a paired positive and negative set of charges appeared on five turtle shells, the
positive and negative charges would usually each be incised five times, either as identical
``carbon copies'' or, usually on the later shells in the set, in abbreviated form.\footnote[51]{For the way in which charges might be abbreviated, see Chang Ping-ch'üan (1956), pp. 246, 253-254; Chou Hung-hsiang (1969), pp. 60-69. For a detailed case study, \ref{sec:chapters-ch03:3.7.1} (see see sec. 3.7.1).I. Chang (Ping-pien 24, k'ao-shih, pp. 51-52) notes an interesting case in which divination topics were recorded on the fifth plastron, but not on the first, even though they were apparently cracked on the first.} 51 No set
appears on more than five scapulas or plastrons.\footnote[52]{Chang Ping-ch'üan (1957), p. 10; 1960), p. 400.} 52 Once a bone or shell had been used for
a certain-numbered crack in a set, it is generally true that all the other sets divined on that
bone or shell employ the same crack number (as in fig. 10, an example from period V)
and possibly other crack-numbers.\footnote[53]{E.g., Ping-pien 12-21; K'u-fang 1595 D). On Ping-pien 311 some of the charges are associated with crack 1 only, some with cracks I and 2, some with cracks 1 to 4, some with cracks I to 5, and some with cracks 1 to 10. The scapula fragment, Menzies 2462, is an apparent exception to this generalization, containing number I cracks for one topic and number 3 cracks for another.}53 The fact that part of a bone or shell might be used for,
say, crack number 1 in a set and then some days later reused for crack number 1 in a different
set suggests that the Shang diviners may, on occasion, have employed a primitive filing
system which enabled them to find a particular bone (shell) or set of bones (shells) as
needed.\footnote[54]{E.g., Ping-pien 34-38 contains two sets of divinations, one divined on chia-ch'en (day 1), the other on yi-mao (day 52); the first plastron (Ping-pien 34) contains the number I cracks for both sets, the second plastron (Ping-pien 35 contains the number 2 cracks for both sets, and so on. (see Chang Ping-ch'üan [1957], p. 10; [1960], p. 400; Ping-pien 34, k'ao-shih, p. 63). Such sequential rigor is particularly noticeable in the period V ten-day divinations; in one case, for example, a set of five scapulas appears to have been reused at ten-day intervals in the same sequence for at least seventy days (Hsü-ts'un 2.972 [D fig. 10]). The three crack numbers still present on the fragment are all number 3; presumably each of the five missing crack numbers was number 3 also. For similar examples, see Hsü-ts'un 1.2580; 1.2684; Yi-chu 246.}54
The routine use of sets suggests that repeated cracking was an important element of
at least some Shang divinations. I suppose that, as with other systems of divination elsewhere,
reading the omens was a sensitive and uncertain task which was made easier by having one
crack to check on another.\footnote[55]{Cf. Gadd (1948), pp. 56-57; Parke and Wormell (1956), pp. 39-40; Oppenheim (1964), pp. 214, 227. 39 1] ] J} 55 As with tossing a coin, not once but many times, the use of
sets presumably gave a certain weight and authority to the eventual reading and also, by
multiplying the number of cracks that were made for any one charge, increased the chance
was not completed on another plastron, there do
seem to have been cases where the diviner gave
up before he had made ten cracks, perhaps due
to lack of response (see n. 48) or perhaps, more
simply, because he ran out of space on the plastron,
e.g., Ping-pien 126.14-17; 189.1-2.
% source: scan 057, printed 40
that some of the cracks would be readable.\footnote[56]{There is some evidence to support this view. Sets which were cracked more times stood a better chance of being inscribed with crack notations. Of the thirteen extra-large, twenty-crack sets listed in n. 47, twelve (92 percent) contain one or more crack notations; of the twenty mediumlarge, eleven-to-nineteen crack sets (n. 48), fourteen (70 percent) contain one or more crack notations. By contrast, only 59 percent of the regular ten-crack sets (n. 46) contain crack notations. The function of sets will be considered more fully in Studies.}\footnote[57]{See, for example,\ref{sec:chapters-ch03:3.7.4.2} (see  sec. 3.7.4.2).}\footnote[58]{The crack notations were never signed. It is possible that the routine crack notations, which make no mention of the king, were determined by the diviner; the more formal prognostications, almost invariably made by the king (sec. 2.7), may have served to confirm the interpretation reached by his ritual assistants. In this view, prognostication would have been the king's prerogative, but he would have been able to delegate the routine reading of individual cracks.}\footnote[59]{Ts'ui-pien 9, k'ao-shih, p. 4; Chang Pingch'üan (1952), p. 621; (1957), p. 6.}\footnote[60]{Cf. Chang Ping-ch'üan (1954), p. 227.}\footnote[61]{Chang Ping-ch'üan (1967), p. 861.}\footnote[62]{Chang Ping-ch'üan (1954), p. 227, citing Ping-pien 96.25. The fact that a crack number sometimes appears in the midst of the crack notation pu tsai ming (?) (\ref{ch:4} (see ch. 4), n. 134) suggests that crack numbers were engraved before crack notations (Huang P'ei-jung [1975], pp. 3a—b).}\footnote[63]{Preliminary research indicates that numberless cracks (cf. n. 32) never had crack notations recorded by them, though they might belong to sets in which crack notations were recorded for another, numbered crack (e.g., Ping-pien 402).}\footnote[64]{The graph hsiung X, the standard Yi-ching notation for ``inauspicious,'' does not appear in the Shang oracle-bone inscriptions (cf. \ref{ch:3} (see ch. 3), n. 89).}
% source: scan 058, printed 41
plastrons, they generally were limited to the back of scapulas.\footnote[65]{Cf. Chou Hung-hsiang (1969), p. 37. In what I have referred to as ``display inscriptions'' (defined in n. 90), the preface, charge, prognostication, and verification were all recorded together on the front of the shell (Keightley [1973], pp. 10-13; [1975], pp. 13-14); e.g., Ping-pien 57.1; 235.2 (fig. 11); 247.1 (fig. 12); Ching-hua 1-6 (S307.3). But in period I at least, the inscription units might also be split in various ways. For example, one of a pair of charges might appear on the front of a shell, the other on the back, as in Ping-pien 65.1/66.1; 65.4/66.4; 65.7/66.9; 276.2/ 277.2. Or the preface and charge (and even verification) might be on the front, with the prognostication on the back, as in Ping-pien 8/9; 28.1/29.1; 28.2/29.2; 61.3-4/62.1. See too \ref{ch:4} (see ch. 4), n. 116. Or, not infrequently, the charge might appear on the front, while the preface, prognostication, and verification would be on the back of the shell, as in Ping-pien 153.13/154.2-3; 212.1/213.1-2; 225/ 226.2 (for these and other examples, see Li Taliang [1972], pp. 69-79).} 65 The standard introductory
formula was wang chan yüeh (S307.3-311.1), ``The king, reading the cracks,
said:\footnote[66]{A fuller form, wang chan pu yüeh, occasionally appears (e.g., Ping-pien 382.3; K'ufang 1935 [D S282.1]); in these cases I take pu to be nominal, ``The king, reading the cracks, said....'' The oracle-bone graph is composed of a pu crack over a mouth, placed within a bone, presumably a scapula. The mouth-plus-crack element may have served as phonetic (GSR 618), with the bone as signific, the graph depicting ``speaking about the crack made in the bone'' (cf. CKWT, p. 1112); the correctness of this view is supported by the fact that the yüeh, essentially redundant, might on occasion be omitted, as in Ping-pien 235.2 (fig. 11). The graph is thought to be ancestor to chan 5, ``prognosticate,'' which in later times lost its pyromantic connotations and bone signific. (The form does appear occasionally in the RFG inscriptions [S506.1], though its meaning is not always certain.) There is no support for the suggestion of Hopkins (1938), p. 416, given currency by its appearance in Soothill (1951, pp. 141-142, that wang in this formula referred to the dead kings.} ... ``66 These forecasts, which were carved for only a minority of the charges,\footnote[67]{We cannot tell whether prognostications were made in every case and recorded in only a few or whether they were made only in those few cases for which they were recorded. Certain divination topics, such as those concerned with the routine ritual schedule, presumably required no formal prognostication (or verification), though on this point, see \ref{ch:4} (see ch. 4), n. 99; the crack notation sufficed. For other topics, it is difficult to determine what percentage of charges was accompanied by prognostications (or verifications); this is because a charge carved ten (five positive, five negative) times might require the recording of only one prognostication (or verification). A ratio of ten charges to one prognostication (or verification), therefore, would indicate that 100 percent (and not 10 percent) of the charges had been prognosticated. These considerations seriously qualify the preliminary discussion in Keightley (1973a), pp. 3–6.} 67
sometimes merely served as a more elaborate crack notation; for example, wang chan yüeh
chi (S308.2), ``The king, reading the cracks, said: 'Auspicious.'''' Sometimes
the prognostication reiterated the wording of the charge; for example, ``The king, reading
the cracks, said: 'Auspicious. We will receive this harvest' '' (fig.9). And, in a certain number
of cases, the terms of the prognostication were more specific than those of the charge itself;
for example, for the bifurcated topic ``Fu Hao's childbearing will be good / Fu Hao's
childbearing will not perhaps be good” we find, ``The king, reading the cracks, said: 'If it
be a ting-day childbearing, it will be good. If it be a keng-day childbearing, it will be
extremely auspicious. If it be a jen-hsü (day 59) (childbearing), it will be inauspicious''''
(fig.\footnote[68]{Ping-pien 248.7. See too the prognostications on Ping-pien 17 (fig. 16) and 19 (sec. 3.7.1.2). The cases in which the prognostication is more specific than the charge provide a significant clue to the way the cracks were interpreted (sec. 3.7.4); the matter will be discussed fully in Studies. 41 U 0 ]} 13).68
This generally describes the situation in period I; for the evolution of the prognosti-
% source: scan 059, printed 42
cations, see sec. 4.3.1.8. The prognostication in any period was usually made by the king;
it was never made by--or, more accurately, the record never attributes it to-one of the
court diviners. They divined, but the king prognosticated. The only other prognosticators
who appear in the record are certain diviners of the Royal Family group.\footnote[69]{The diviners and Hsieh \#+ prognosticated (Ch'ien-pien 4.42.2 [S318.2]; 6.39.1; Chingchin 1599; 1601 [all S311.1]; Hsü-pien 5.7.4 [S506.1]; Wai-pien 240 [S13.3]; Yi-ts'un 586 [S139.2, but wrongly recorded as 584]; Yin-hsü 574 [D S311.1]). These few prognostications, written in small script, usually containing just one word, chi, ``auspicious,'' and frequently omitting the chan, may be regarded as different in kind from the boldly written and, on occasion, remarkably precise prognostications of the period I court divinations.}\footnote[70]{E.g., Ching-hua 1-6 (S307.3.}\footnote[71]{Ping-pien 248.7 (fig. 13).}\footnote[72]{Yi-pien 6668/6669 (S240.1. Jao (1961), p. 971, n. 2, concerned by the long interval, suggests that the prognostication was made two days before the charge was divined!}\footnote[73]{For a preliminary study, see Keightley (1975c), pp. 11–17.}\footnote[74]{For the difficulties involved in estimating the ratio of charges, prognostications, and verifications, see n. 67.}\footnote[75]{See n. 65.}
\inscriptionsection{PAIR OF COMPLEMENTARY CHARGES}

\begin{inscriptionblock}
\inscriptionside{positive, left side}
\inscriptionref{Ping-pien 235.2}
Crack-making on chi-mao (day 16),
Ch'üch divined:
``It will rain.''
The king, reading the cracks (said):
``It will rain; it will be a jen day.''
On jen-wu (day 19) it really did rain.

\inscriptionside{negative, right side}
\inscriptionref{Ping-pien 235.1}
Crack-making on chi-mao (day 16),
Ch'üch divined:
``It will not perhaps rain.”

\end{inscriptionblock}

In period I, verifications (and prognostications) might provide considerable detail; for
example:

\inscriptionsection{PAIR OF COMPLEMENTARY CHARGES}

\begin{inscriptionblock}

\inscriptionside{negative, left side}
\inscriptionref{Ping-pien 1.4}
Crack-making on kuei-ch'ou (day 50),
Cheng divined:
``From today to ting-ssu (day 54) we
will not perhaps harm the Chou.''
\inscriptionside{positive, right side}
\inscriptionref{Ping-pien 1.3}
Crack-making on kuei-ch'ou (day 50),
Cheng divined:
``From today to ting-ssu (day 54) we
will harm the Chou.''
The king, reading the cracks, said:
``(Down to) ting-ssu (day 54) we
should not perhaps harm (them);
on the coming chia-tzu (day 1) we
will harm (them).''

\end{inscriptionblock}

On the eleventh day,\footnote[76]{On the possible functions of wu, see Takashima (1973), pp. 91-114. Cf. \ref{ch:3} (see ch. 3), n. 63. On ch'i, see \ref{ch:3} (see ch. 3), n. 41.}\footnote[77]{Ordinal day counts in period I inscriptions run inclusively from the day of divination. It is possible that RFG counts excluded the day of divination, but the evidence is uncertain. The issue will be discussed in Studies.}\footnote[78]{Chang Ping-ch'üan (Ping-pien 1, k'ao-shih, p. 6) takes ch'e as a name; Takashima (1973), p. 110, takes it as a chariot unit. Either interpretation is possible.}\footnote[79]{The graph (which, for the sake of convenience, I romanize as t'ou = [?], a vessel without a cover [CKWT, p. 4070]) appears in two different contexts in the inscriptions: \listitem{1} as a verb of sacrifice (S403.1); see Chang Tsungtung (1970), p. 71, n. 5; Serruys (1974), p. 70; \listitem{2} as a link between two kan-chih dates which are always consecutive, the formula being chia-tzu (hsi) t'ou yi-ch'ou F \listitem{5} Z (S402.3-4); see Ping-pien 96, k'ao-shih, pp. 135-137. In this second case, t'ou apparently referred to the nighttime no-man's-land between two cyclical-day dates, and that is the usage employed here. For a preliminary analysis, see Keightley (1975), pp. 142144; a fuller discussion will be provided in Studies. L ]}
% source: scan 061, printed 44
Or, for still more detail: (Ching-hua 2 [fig. 14], left side [S401.4]):

\inscriptionsection{INSCRIPTION}

\begin{inscriptionblock}
\inscriptionref{Ching-hua 2 [fig. 14], left side [S401.4]}

(PREFACE:) Crack-making on kuei-ssu (day 30), Ch'üeh divined:
(CHARGE:) ``In the (next) ten days there will be no disaster.''
(PROGNOSTICATION:) The king, reading the cracks, said: ``There will be harm; there
will perhaps be the coming of alarming news.'
9980
(VERIFICATION:) When it came to the fifth day,\footnote[80]{"Disasters,'' ``alarming news,” and “harm'' are functional translations (sec. 3.5) of (a) F, (b) and (c), respectively. The Shang apparently envisioned a hierarchy of disaster words which were general or specific, potential or actual. According to Mickel (1976), pp. 79-94, 100-110, 240—who notes the existence of eighteen different disaster graphs from period I-form a was the most generalized term, used vaguely of any disaster that might occur or be forecast in the future; form c, by contrast, referred to disasters that had been forecast or had already occurred. For the functional translation, ``alarming news,'' for form b, cf. Chow Tse-tsung (1968), p. 168, n. 65; Mickel (1976), pp. 174-178.}\footnote[81]{See n. 77.}\footnote[82]{See \ref{ch:3} (see ch. 3), n. 104.}\footnote[83]{For the implications of this point, see Keightley (1975c), pp. 13-15. The prognostication was at times unsupported by the verification, e.g., Ping-pien 368/369; Yi-pien 6419/6420 (depending on the meaning of and, since the rubbing is hard to decipher, following the reading of Li Yen [1969], p. 195); Cho-ts’un 34 (S281.4 = Chui-hsin 530); in these cases, the thing forecast happened, but not strictly in accordance with the terms of the prognostication. I know of only two cases where the recorded prognostication may have been contradicted by the verification: Yi-pien 4130 (S245.4) (if, with Jao [1957/58], p. 18, means ``lose life''); Hsü-ts'un 2.76 (D) (S233.1) (if means ``foggy'').}
\end{inscriptionblock}

% source: scan 062, printed 45
followed fifty-seven and twenty-seven days after his sickness had been prognosticated;
other verifications \listitem{8} involve an interval of a hundred or more days; and in a final case \listitem{9}
the onset of sickness appears to have taken place 175 (?) days after the divination.84
In this last case, as in the others, the oracle bone may have been kept on file, with the
verification (and perhaps the prognostication and charge as well)\footnote[85]{Cf. n. 90.} 85 uncarved, awaiting
the final confirmation, or confutation, of the royal forecast.

\section[Crack Notations]{\origsecnum{2.6}Crack Notations}
\label{sec:chapters-ch02:2.6}
that some of the cracks would be readable.There is some evidence to support this view. Sets which were cracked more times stood a better chance of being inscribed with crack notations. Of the thirteen extra-large, twenty-crack sets listed in n. 47, twelve (92 percent) contain one or more crack notations; of the twenty mediumlarge, eleven-to-nineteen crack sets (n. 48), fourteen (70 percent) contain one or more crack notations. By contrast, only 59 percent of the regular ten-crack sets (n. 46) contain crack notations. The function of sets will be considered more fully in Studies.See, for example,\ref{sec:chapters-ch03:3.7.4.2} (see  sec. 3.7.4.2).The crack notations were never signed. It is possible that the routine crack notations, which make no mention of the king, were determined by the diviner; the more formal prognostications, almost invariably made by the king (sec. 2.7), may have served to confirm the interpretation reached by his ritual assistants. In this view, prognostication would have been the king's prerogative, but he would have been able to delegate the routine reading of individual cracks.Ts'ui-pien 9, k'ao-shih, p. 4; Chang Pingch'üan (1952), p. 621; (1957), p. 6.Cf. Chang Ping-ch'üan (1954), p. 227.Chang Ping-ch'üan (1967), p. 861.Chang Ping-ch'üan (1954), p. 227, citing Ping-pien 96.25. The fact that a crack number sometimes appears in the midst of the crack notation pu tsai ming (?) (\ref{ch:4} (see ch. 4), n. 134) suggests that crack numbers were engraved before crack notations (Huang P'ei-jung [1975], pp. 3a—b).Preliminary research indicates that numberless cracks (cf. n. 32) never had crack notations recorded by them, though they might belong to sets in which crack notations were recorded for another, numbered crack (e.g., Ping-pien 402).The graph hsiung X, the standard Yi-ching notation for ``inauspicious,'' does not appear in the Shang oracle-bone inscriptions (cf. \ref{ch:3} (see ch. 3), n. 89).

\section[The Prognostication]{\origsecnum{2.7}The Prognostication}
\label{sec:chapters-ch02:2.7}
plastrons, they generally were limited to the back of scapulas.Cf. Chou Hung-hsiang (1969), p. 37. In what I have referred to as ``display inscriptions'' (defined in n. 90), the preface, charge, prognostication, and verification were all recorded together on the front of the shell (Keightley [1973], pp. 10-13; [1975], pp. 13-14); e.g., Ping-pien 57.1; 235.2 (fig. 11); 247.1 (fig. 12); Ching-hua 1-6 (S307.3). But in period I at least, the inscription units might also be split in various ways. For example, one of a pair of charges might appear on the front of a shell, the other on the back, as in Ping-pien 65.1/66.1; 65.4/66.4; 65.7/66.9; 276.2/ 277.2. Or the preface and charge (and even verification) might be on the front, with the prognostication on the back, as in Ping-pien 8/9; 28.1/29.1; 28.2/29.2; 61.3-4/62.1. See too \ref{ch:4} (see ch. 4), n. 116. Or, not infrequently, the charge might appear on the front, while the preface, prognostication, and verification would be on the back of the shell, as in Ping-pien 153.13/154.2-3; 212.1/213.1-2; 225/ 226.2 (for these and other examples, see Li Taliang [1972], pp. 69-79). 65 The standard introductory
formula was wang chan yüeh (S307.3-311.1), ``The king, reading the cracks,
said:A fuller form, wang chan pu yüeh, occasionally appears (e.g., Ping-pien 382.3; K'ufang 1935 [D S282.1]); in these cases I take pu to be nominal, ``The king, reading the cracks, said....'' The oracle-bone graph is composed of a pu crack over a mouth, placed within a bone, presumably a scapula. The mouth-plus-crack element may have served as phonetic (GSR 618), with the bone as signific, the graph depicting ``speaking about the crack made in the bone'' (cf. CKWT, p. 1112); the correctness of this view is supported by the fact that the yüeh, essentially redundant, might on occasion be omitted, as in Ping-pien 235.2 (fig. 11). The graph is thought to be ancestor to chan 5, ``prognosticate,'' which in later times lost its pyromantic connotations and bone signific. (The form does appear occasionally in the RFG inscriptions [S506.1], though its meaning is not always certain.) There is no support for the suggestion of Hopkins (1938), p. 416, given currency by its appearance in Soothill (1951, pp. 141-142, that wang in this formula referred to the dead kings. ... ``66 These forecasts, which were carved for only a minority of the charges,We cannot tell whether prognostications were made in every case and recorded in only a few or whether they were made only in those few cases for which they were recorded. Certain divination topics, such as those concerned with the routine ritual schedule, presumably required no formal prognostication (or verification), though on this point, see \ref{ch:4} (see ch. 4), n. 99; the crack notation sufficed. For other topics, it is difficult to determine what percentage of charges was accompanied by prognostications (or verifications); this is because a charge carved ten (five positive, five negative) times might require the recording of only one prognostication (or verification). A ratio of ten charges to one prognostication (or verification), therefore, would indicate that 100 percent (and not 10 percent) of the charges had been prognosticated. These considerations seriously qualify the preliminary discussion in Keightley (1973a), pp. 3–6. 67
sometimes merely served as a more elaborate crack notation; for example, wang chan yüeh
chi (S308.2), ``The king, reading the cracks, said: 'Auspicious.'''' Sometimes
the prognostication reiterated the wording of the charge; for example, ``The king, reading
the cracks, said: 'Auspicious. We will receive this harvest' '' (fig.9). And, in a certain number
of cases, the terms of the prognostication were more specific than those of the charge itself;
for example, for the bifurcated topic ``Fu Hao's childbearing will be good / Fu Hao's
childbearing will not perhaps be good” we find, ``The king, reading the cracks, said: 'If it
be a ting-day childbearing, it will be good. If it be a keng-day childbearing, it will be
extremely auspicious. If it be a jen-hsü (day 59) (childbearing), it will be inauspicious''''
(fig.Ping-pien 248.7. See too the prognostications on Ping-pien 17 (fig. 16) and 19 (sec. 3.7.1.2). The cases in which the prognostication is more specific than the charge provide a significant clue to the way the cracks were interpreted (sec. 3.7.4); the matter will be discussed fully in Studies. 41 U 0 ] 13).68
This generally describes the situation in period I; for the evolution of the prognosti-

\section[The Verification]{\origsecnum{2.8}The Verification}
\label{sec:chapters-ch02:2.8}
cations, see sec. 4.3.1.8. The prognostication in any period was usually made by the king;
it was never made by--or, more accurately, the record never attributes it to-one of the
court diviners. They divined, but the king prognosticated. The only other prognosticators
who appear in the record are certain diviners of the Royal Family group.The diviners and Hsieh #+ prognosticated (Ch'ien-pien 4.42.2 [S318.2]; 6.39.1; Chingchin 1599; 1601 [all S311.1]; Hsü-pien 5.7.4 [S506.1]; Wai-pien 240 [S13.3]; Yi-ts'un 586 [S139.2, but wrongly recorded as 584]; Yin-hsü 574 [D S311.1]). These few prognostications, written in small script, usually containing just one word, chi, ``auspicious,'' and frequently omitting the chan, may be regarded as different in kind from the boldly written and, on occasion, remarkably precise prognostications of the period I court divinations.E.g., Ching-hua 1-6 (S307.3.Ping-pien 248.7 (fig. 13).Yi-pien 6668/6669 (S240.1. Jao (1961), p. 971, n. 2, concerned by the long interval, suggests that the prognostication was made two days before the charge was divined!For a preliminary study, see Keightley (1975c), pp. 11–17.For the difficulties involved in estimating the ratio of charges, prognostications, and verifications, see n. 67.See n. 65.

\section[Recording the Inscriptions]{\origsecnum{2.9}Recording the Inscriptions}
\label{sec:chapters-ch02:2.9}

Having dealt first with the form and content of the inscriptions, I want to turn
now to the mechanics of oracular record keeping. Contrary to what some scholars have
supposed,\footnote[86]{E.g., Britton (1940), pp. 7-8; Thompson (1969), p. 34; Wilkinson (1973), p. 60.}86 the engravers did not carve the charge into the bone or shell before the cracking
took place, but after.\footnote[87]{Ch'en Meng-chia (1956), p. 13; Chang Ping-ch'üan (1967), p. 866. That the inscriptions were generally made after the cracks is shown by the fact that in only a few instances is a crack straddled by a graph written across it, a situation known as fan-chao (sec. 2.4); for examples, see Ping-pien 1; 96 (esp. 96.25); 155; 159; 243; 284; 502; Yi-pien 867. Had the cracks not been formed until after the engravers had carved the graphs, we would expect many more cases of fan-chao; it would have been extremely tedious, if not impossible, for the carver to keep turning the bone or shell over to estimate on the basis of the hollow locations on the back the places on the front where he should not write for fear that cracks would later bisect the graph. That the inscriptions were carved after the cracks is also shown by the following considerations. If we agree that the crack numbers (sec. 2.4) were carved after the diviners had made the cracks (it is hardly likely that the crack numbers would have been carved before) and if we agree (on the basis of crack numbers being erased to make way for inscriptions; see n. 39) that the inscriptions were carved after the crack numbers, then it follows that the inscriptions must have been carved after the cracks were made.} 87 This suggests that the scapulimantic and plastromantic inscriptions were not simply regarded as prayers or magical letters forwarded to the spirits. The
writing recorded not what was about to be divined but what had been. Its purposes, therefore, were at least partly historical and bureaucratic—to identify the topics, forecasts, and
results for which the cracks had been formed.\footnote[88]{No rigorous study has been made of which cracks on a shell were associated with which inscriptions; there is no way to be certain, for example, that the Ping-pien k'ao-shih assigns the right cracks to the right inscriptions in every case. Quite probably, the Shang diviners observed no very strict rules themselves, though proximity must have been the governing factor and boundary lines (sec. 2.11) were, on occasion, carved to avoid confusion. On this point, cf. Chang Ping-ch'üan (1956), p. 239; see too\ref{sec:chapters-ch03:3.7.2} (see  sec. 3.7.2). 45 Q D E L U [}88
How soon after the cracks were burned the divination record was carved we cannot
84. \listitem{1} Ching-hua 1 (\$307.3); cf. Shirakawa
(1972), p. 26; \listitem{2} Ping-pien 102.1; \listitem{3} Ping-pien
235.2 (fig. 11); \listitem{4} Ching-hua 2 (S401.4;
% out-of-scope-heading: 154.3 235.2 (fig. 11); \listitem{4} Ching-hua 2 (S401.4;

translated in\ref{sec:chapters-ch02:2.8} (see  sec. 2.8)); \listitem{5} Yi-chu 620 (S91.1);
\listitem{6} Ping-pien 247.1-2 (fig. 12); \listitem{7} Yi-pien 4130;
(S245.4); it is not certain, however, that Jao@@
(1957/58), p. 18, is correct in taking the second
inscription which certainly does involve a verification fifty-seven days later—as referring to the
death (?) of the same person; \listitem{8} Pu-tz'u 579 and
other inscriptions at S467.2; \listitem{9} T'ieh-yün 5.3
(S243.4); various reconstructions of certain of the
cyclical dates are possible (e.g., Tung-tsuan 530,
k’ao-shih, p. 116b; Yen Yi-p'ing [1952], p. 296),
but it seems that an interval of 170-odd days was
involved between divination (or prognostication)
and verification.
% source: scan 063, printed 46
tell.89 And it is not certain whether all the units of an inscription\footnote[89]{I do not find persuasive the arguments of Hu Hou-hsüan (1939), pp. 412-413, 425, which attempt to show, on the basis of such matters as omissions in the record, that the charges were not necessarily incised on the same day. His evidence may be explained by intentional abbreviation of the record, subsequent fragmentation of the bones, etc., rather than by the forgetfulness he imputes to the recorders.} were carved at the same
time or whether the various units were inscribed separately, with the preface and charge
being inscribed first, for example, the prognostication being inscribed after the cracks had
been read, and the verification being inscribed only after events had taken their course.\footnote[90]{The argument that the prognostications and verifications were inscribed significantly later than the charges (e.g., Britton [1940], p. 8; Tung [1954b], p. 359; Li Ta-liang [1972], p. 109) depends either upon the claim that this would have been logical or that the writing style of these units differs from that of the rest of the inscription. But what is logical to us may not have been logical to the Shang. And it is frequently hard to establish the existence of a significant change in the writing style (e.g., in Ping-pien 32.27; 235.2; 251.2; 334.3, are the graphs of the prognostication larger, or smaller, than those of the charge?); and even when the change is evident with the start of the verification (as in Ping-pien 284.4), we cannot be certain that the smaller size and slightly different angle with which the graphs of the verification were written were not signs of lack of space rather than of a later record. On the other hand, certain ``display inscriptions'' appear to have been written all at one time. As I define them, the essential characteristics of these display inscriptions are: \listitem{1} bold, large calligraphy; \listitem{2} the prognostication and verification written as a single, continuous unit, and usually placed immediately next to the charge; \listitem{3} the verification, frequently detailed, confirms the accuracy of the prognostication (Keightley [1975], p. 13). For typical display inscriptions, see Ping-pien 1.3 (translated in\ref{sec:chapters-ch02:2.8} (see  sec. 2.8)); 57.1; 247 (fig. 12); Ching-hua 2 (S401.4; translated in\ref{sec:chapters-ch02:2.8} (see  sec. 2.8)); K'u-fang 1595 (D) (S301.4302.1); T'ieh-yün 5.3 (S243.4; discussed in n. 84 above). Objective criteria need to be established for determining what was, and was not, written at one time; it is possible, for example, that microscopic studies of engraving strokes (cf. Chou Hung-hsiang [1967/68]) may eventually tell us if, in such cases, the same or different knife and hand carved the whole inscription.}90
When the diviners used a bone or shell over several days, it is probable that the divination
inscriptions of the first day were engraved before the subsequent divinations were made.\footnote[91]{See \ref{ch:3} (see ch. 3), n. 113, for evidence; cf. n. 140.}91
But practice may have varied on all these points, and further research is needed.

\subsection{\origsecnum{2.9.1}Writing Brush}
\label{sec:chapters-ch02:2.9.1}

A small number of the inscriptions were apparently written first with a brush
dipped in ink and were then incised; these graphs are generally large and crude in style, and
are found on the back of scapulas or plastrons.\footnote[92]{Ch’en Meng-chia (1956), p. 14; Chang Ping-ch'üan (1954), p. 231; (1967), p. 866; Chou Hung-hsiang (1967/68), p. 10. See, for example, Yi-pien 6423 (after Tung [1948a], p. 124); KK (1975.1), p. 39, citing H103: 62.} 92 Some marginal notations were also written
on the bridge with a brush.\footnote[93]{Ch’en Meng-chia (1956), p. 13. Ching-chin 2, a plastron back, has a three-character marginal notation written in red, but not carved (Yen Yi-p'ing [1973], p. 2a).} 93 The brush strokes generally followed the same direction and
sequence as those used by traditional and modern calligraphers.\footnote[94]{Tung (1948a), p. 127, and accompanying Chinese text; (1949a), p. 264.} 94 The writing brush was
apparently pliable and springy, and since its strokes were thicker than those subsequently
carved by the knife, traces of the extra ink remain in certain cases.\footnote[95]{E.g., Ping-pien 66.3, 5, 7. On the nature of early writing materials in general, see Su Ying-hui (1957).} 95 We also find instances
% source: scan 064, printed 47
in which inscriptions were first written by brush but, due presumably to clerical oversight
or to the fact that the writing was only for practice, were only partially carved, if at all.\footnote[96]{E.g., Chia-pien 870; 2636; 2940; Ping-pien 66.1, 2, 4, 5, 6. Tung (1965), plates 40-42 (the last originally published in Tung [1933], facing p. 418), gives other photographs of brush-written characters on bone. Such characters do not reproduce in rubbings; cf., for example, the rubbing and photograph provided for Ping-pien 27.2, 4, 5.}96
The ink, originally red (and perhaps black), has with age turned brown, or black with
yellow and brown tinges, or light yellow.\footnote[97]{Chu (1935), ch. 6, p. 17a; Ch'en Meng-chia (1956), p. 14.}97

\subsection{\origsecnum{2.9.2}Engraving Knife}
\label{sec:chapters-ch02:2.9.2}

It appears that the great majority of the inscriptions were carved directly into the
surface of the bone, with no prior brush-written characters to copy.\footnote[98]{Ch'en Meng-chia (1956), p. 15, gives the following reasons for rejecting the theory (e.g., Tung [1949a], p. 263) that the carved graphs required brush-written ``templates”: \listitem{1} the front of the plastron was too slick to take ink readily; \listitem{2} since the brush-written graphs are larger and cruder than the carved ones and since they are also frequently upside down (cf. n. 111 below), they were clearly not templates that had not been carved through an oversight; \listitem{3} some engraved graphs were so small that the engraver could not have followed brush strokes accurately; \listitem{4} the number of commonly used graphs was not large [cf. \ref{ch:3} (see ch. 3), n. 7]; there was thus no necessity for templates to be provided. Ch'en's third reason contradicts the view of Kaizuka and Itō (1953), pp. 74, 76, that the fine, small calligraphy of the Tzu group (sec. 2.2.1.1) and period V inscriptions was first written with a brush; their view is weakened by Ch'en's second point, i.e., the size and crudity of the brush-written graphs we know (as in Ping-pien 66.1).} 98 Many practice
inscriptions have been found,\footnote[99]{The following types of practice inscription, identified by crudity of form, lack of sense, or lack of connection with divinatory cracks (criteria, incidentally, that may also be used to identify fake inscriptions; \ref{sec:chapters-ch05:5.4} (see see sec. 5.4)), may be noted: (I) practice inscriptions carved on unprepared or unburned bone (e.g., KK (1975.1], p. 45, citing H99:3); \listitem{2} practice inscriptions involving senseless repetition of graphs (e.g., Chia-pien 2875; T'ai-ta 2, 7; see Liu Yüan-lin [1974], p. 121, for another case, of uncertain provenance); \listitem{3} practice graphs accompanied by doodles of animal figures (e.g., Chia-pien 2336; 2422; see Hsü Chungshu [1931], p. 528, for a drawing); \listitem{4} practice graphs cut into a bone already used for regular inscriptions (e.g., Chia-pien 2407.3; Chia-t'u 125/ 126 \listitem{5} practice graphs cut over old inscriptions that had been erased (e.g., KK [1975.1], pp. 38-39, citing H85:127); \listitem{6} kan-chih tables (n. 100.}99 and Kuo Mo-jo has suggested that since freehand carving
would have taken considerable practice, the numerous kan-chih day tables which we find
engraved on bone should be regarded as “exercise sheets.\footnote[100]{Kuo detects in one fragment a row of graphs written by the teacher; the other rows are written by his students, whose better efforts were due to the teacher's guiding hand (Ts'ui-pien 1468, k'ao-shih, p. 196b; cf. Kuo Mo-jo [1972a], pp. 4-5). For representative kan-chih tables, see Chia-pien 2599; KK (1975.1), pp. 38-39, citing H85: 127; Chui-hsin 373-376; 415; 433; 434; 475; 679; 682; Hou-pien 2.1.5 (with many horizontal strokes missing; \ref{sec:chapters-ch02:2.9.3} (see see sec. 2.9.3)); US 1; 40 (Chou Hunghsiang [1976], pp. 7, 8).}"100
It is generally thought that the inscriptions were carved with engraving knives;\footnote[101]{Chang Ping-ch'üan (1954), p. 231; Ch'en Meng-chia (1956), p. 13. 47 C D J 1 U}101
they were presumably made of bronze, though some may have been of jade. No Shang
knives have yet been excavated which can be said, with certainty, to have been used for
% source: scan 065, printed 48
carving graphs on bone and shell. 102 We are\footnote[102]{Chou Hung-hsiang (1967/68), on the basis of classical descriptions, other Shang knives, and a study of the grooves left in the Shang bones, reconstructs four kinds of blades and suggests the ways in which the engraving knife may have been held. See too Hsü Chin-hsiung (1974), p. 283, n. 56. A bronze knife, similar to that used by modern engravers, was apparently found in 1929 with the large plastrons excavated from the Talien-k'eng (hereafter TLK) (Matsumaru [1959], p. 8, citing an unspecified report of Tung). Kuo Pao-chün (1951), pp. 26-27, plate 9.4, reported a jade knife, excavated from a Wu-kuants'un grave in 1950, that was still sharp enough to cut strokes on turtle shell (the knife is also reproduced by Itō [1958], fig. 132). For other references to Shang engraving knives, see: Matsumaru (1959), p. 8, who, like Britton (1940), p. 3, argues for a blade with a V-shaped cross-section; Watson (1963), p. 42. Tung (1965), plate 39, provides an enlarged photograph of the grooves forming one graph on Chui-ho 36 (Tung's plate 38); a long thin groove extending from the graph may indicate where the knife slipped. As Chou Hunghsiang (1967/68), p. 18, points out, it is hard to cut graphs on bone or shell today, even with a modern steel blade; it would have been far harder with a jade blade. Britton (loc. cit.) put it well: ``The osteoglyphs of the Ti Yi and Ti Hsin reigns [i.e., period V] are the despair of present-day imitators. One marvels at the strength and control of wrist and fingers which the writers must have developed.'' Barnard (1960a), p. 482, found that ``a gentle scraping action rather than deliberate cutting seems to be the more reliable way to achieve clear-cut strokes,'' but, concerned with the way forgeries are made, he was working with old bone from a Shang site, not, as the Shang diviners were, with fresh bone. Chou Hung-hsiang (1976), p. 12, has noted one piece, US 486, at Harvard University, with ``inscriptions engraved by a gouge rather than a knife."} still not sure, in fact, how the Shang, with
their comparatively primitive blades, were routinely able to cut characters into the hard
bone and shell that were regularly and neatly formed and, as in period V (for example,
fig. 10), small and remarkably finely shaped.\footnote[103]{Kuo Mo-jo (1972a), p. 4, has suggested that by analogy with ivory working, the bone must have been softened in an acid bath; but I know of no evidence that the Shang used acid. Chou Hung-hsiang (conversation, 10 September 1970) proposes that the Shang may have softened the bone or shell by boiling it.} 103

\subsection[Engravers]{\origsecnum{2.9.3}Engravers}
\label{sec:chapters-ch02:2.9.3}

The writing style associated with one diviner might vary greatly,\footnote[104]{Jao (1959), pp. 1188-1189, lists many cases of varied calligraphic styles associated with single diviners; cf. Barnard (1960a), p. 485. In period I, to cite a few examples, we find the divinations of Cheng recorded in both large, copperplate (Pingpien 1.1-4) and medium-sized calligraphy (Pingpien 3.1-2; 5.1-2). The divinations of Ch'üeh were likewise recorded in large (Ping-pien 8 [fig. 8]), medium (Ping-pien 2.1-2; 12; 14; 16 [fig. 15]; 18; 20.3-6), and small (Ping-pien 1.9-18; 2; 3.21) calligraphy. In period II, the calligraphy of diviner Yi's inscriptions might be either medium (Ch'ien-pien 6.21.2; Jimbun 1588) or small (Houpien 2.13.2; 2.24.4); the diviner Chung + 's inscriptions (S570-571) also varied in style; on the other hand, the inscriptions of the diviner Hsing (S567-568) were invariably small (on the meaning of the terms ``copperplate,” “large,'' ``medium,'' etc., see sec. 4.3.1.3). For epigraphic variations in the names of particular diviners, see, for example, Ch'ü Wan-li (1960a), p. 189, and Li Yen (1966/67) on the diviner Ho J (cf. Cheng Te-k'un [1971], p. 146; Barnard [1973], p. 24); Yen Yi-p'ing (1961), p. 506, on the diviners and Wei Tung (1936), pp. 107-113, on the diviner ýt (cf. Ch’en Meng-chia [1956], p. 16).} 104 yet the writing
style associated with different diviners was frequently identical;\footnote[105]{The engraving style of certain period I diviners like Cheng, Ch'üeh, Hsüan, and Pin is so similar that it is reasonable to assume they were the product of the same hand. Indeed, it is the similarity of style for diviner groups that permits calligraphy and epigraphy to be used as a dating criterion (secs. 4.3.1.3 and 4.3.1.4). For a good example of identical writing style used for the inscriptions of two RFG diviners, see Ping-pien 612, k'ao-shih, pp. 94-95.} 105 and one diviner's
% source: scan 066, printed 49
inscriptions and even the form of his name might be written differently on the same
plastron.\footnote[106]{Ping-pien 1.1; 1.3; and 1.18 all record divinations by the diviner Cheng over a fourteenday period; the engraving style of each is markedly different. Similarly, the divinations of the diviner Ch'üeh appear in two different styles on the Ping-pien 12-21 set (sec. 3.7). Chia-t'u 113 records two divinations by Ho, made ten days apart; the epigraphy differs markedly for a number of the graphs. For other instances, see the sources cited in n. 104. On at least one occasion, the engraver changed in mid-inscription; the k'ao-shih to Chiapien 2779 notes that the last five graphs of the charge were cut by a different hand.} 106 These facts indicate that the inscriptions were not engraved by the diviners
themselves, but by special engravers,\footnote[107]{Ch'en Meng-chia (1956), p. 16, commenting on Tung (1936), pp. 107-113; see too Barnard (1973), p. 26, n. 8.} 107 presumably those who were being trained to write
the kan-chih tables referred to in\ref{sec:chapters-ch02:2.9.2} (see  sec. 2.9.2). The hypothesis of Matsumaru that the records
were incised by a comparatively small number of engravers-reduced perhaps to only
one man in period V, when the calligraphy is remarkably uniform—accords with the
evidence,\footnote[108]{Matsumaru (1959), p. 8. An estimate of the number of bones and shells engraved per day (appendix 3, sec. 4) supports the view that the engravers could have been few in number. The existence of individual engravers, whose careers were not necessarily coterminous with those of the diviners or diviner groups, bears on the criteria used for relative dating (secs. 4.3.1.3 and 4.3.1.4).} 108
At least some of the carving appears to have been performed on a mass-production
basis in the sense that the graphs were not all carved one at a time. In these cases,
the engraver first cut all the vertical and sloping strokes for every graph on the bone and then
rotated the bone to carve all the horizontal strokes. Thus, the graphs ting-hai T, whose
full oracle-bone form was , would first be carved as\footnote[109]{Tung (1929), p. 75.} .109 And the carved strokes
gibberish in themselves, may be completed to read\footnote[110]{Hou-pien 1.16.6 (following the reconstruction of Ikeda [1964]; for two earlier reconstructions that differ, see Yeh Yü-sen [1929], p. 10a, and Tung [1933], p. 418). By the nature of the evidence we only know that graphs were carved in this way when strokes were omitted. Nevertheless, the frequent absence of horizontal strokes suggests that despite the considerable experience that must have been required to engrave long inscriptions in first one dimension and then the othermany inscriptions were carved in this way. There are, for example, numerous examples of oracle bones in which a series of horizontal strokes have not been carved. Chou Hung-hsiang (1976), p. 1 I, cites USB 416, a scapula at Columbia University, in which a series of horizontal strokes are missing in the right central section of the bone. Hou-pien 2.1.5 is a notable case in which the graphs in the first two right-hand columns were completely carved, but the horizontal strokes for the graphs in the next six columns are, with one exception, missing (cf. Kuo Mo-jo [1972a], p. 4). For other examples, see Chia-pien 221, k'ao-shih, p. 35; Pingpien 18.7; 30.4; 34-3-4; 37.1, 3, 4 (after Chang Ping-ch'üan [1956], p. 255); Yi-pien 3389.5; Chien-shou 46.14 (after Tung [1933], p. 418); Tung (1936), p. 114; Ch’en Meng-chia (1956), p. I3.} .110 In this way,
the engraver could have plied his knife in one general axis and would have needed to rotate
the bone only once, thus saving time. This method of carving undoubtedly contributed
to the not infrequent errors, erasures, and corrections found in the inscriptions.\footnote[111]{The following types of errors may be noted (the list of categories is taken from Hu Hou-hsüan [1939], pp. 401-421, and Ch'ü Wan-li [1961], p. 5, who give numerous examples-not all of which were necessarily the Shang errors these authors take them to be): \listitem{1} graphs omitted (e.g., Chia-t'u 80.17; but many cases may represent intentional abbreviation;\ref{sec:chapters-ch03:3.7.1.1} (see  cf. sec. 3.7.1.1)); \listitem{2} surplus graphs (e.g., Chia-pien 756.2); \listitem{3} graphs inserted (e.g., Chia-pien 27); \listitem{4} wrong graphs, uncorrected (e.g., Chia-pien 476.1; 2511; Ping-pien 73.6); \listitem{5} wrong graphs, corrected (e.g., Chia-pien 2214); \listitem{6} strokes missing (see n. 110); \listitem{7} graphs upside down (e.g., Chia-pien 1510; see too Yen 49 1 U U}111
% source: scan 067, printed 50

\subsection*{Inscription Placement}

The engravers generally carved the oracle inscriptions on the front of the scapulas
and plastrons, placing them in crack-free areas or fitting them around the cracks. As a rule,
the inscriptions appear to have been carved above, or to the side of, their pu cracks and on
the side of the crack which lacked the transverse branch (for example, fig. 8); but the
modern scholar cannot always be certain which inscription referred to which crack(s).\footnote[112]{See n. 88.} 112
Inscriptions were rarely carved on the back of the bones or shells, except in the reign of
Wu Ting when, as we have seen, units of the same inscription might appear on both the
front and back of a shell.\footnote[113]{See n. 65. For the meaning of ``front'' and ``back,'' see \ref{ch:1} (see ch. 1), n. 46.} 113 The significance of this is uncertain, but it reminds us that
the context of an inscription, especially during period I, involves both the back and front
of the bone or shell on which it appears.\footnote[114]{Tung (1945), pt. 2, ch. 9, p. 1b (cf. Li Ta-liang [1972], p. 47), suggests that some inscriptions ran over to the back because of the lack of crack-free space on the front.} 114
The individual charges on both scapulas and plastrons were generally carved in vertical
columns, and were read downward like traditional Chinese texts.\footnote[115]{For the meaning of ``top'' and ``bottom,'' see \ref{ch:1} (see ch. 1). n. 46. There are a few curious exceptions to the standard orientation. For example, the front of Chia-pien 2902/2903 has to be read with the socket at the bottom, but the back must be read with the socket at the top (this scapula is reproduced by Umehara [1964], plate 23). Some of the inscriptions on Yi-pien 407, a half carapace, are upside down (Tung [1962]). Kuo Pao-chün (1951), plate 41.1-2, reproduces a scapula in which the socket has to be placed at the right in order for the writing to be read. 394.} 115 In nonoracular
Shang writing, the columns normally moved from right to left,\footnote[116]{Tung (1933), pp. 420-421; (1951a), p.} 116 but the natural symmetry
of the plastrons (and scapula pairs)\footnote[117]{Cf. n. 45.} 117 encouraged the oracle-bone engravers to balance
the inscriptions in antipodal fashion so that if one column moved from right to left, the
matching column would move from left to right.
On turtle shells, the columns might be read moving from the center of the plastron
toward the edge, or from the edge toward the center; pairs were usually of this form,
being carved so that the inscriptions either moved toward, or away from, each other, like
mirror reflections (fig.\footnote[118]{Ping-pien 69.1-2; 354.1-2.} 8).118 Charges were occasionally written as single horizontal
lines;\footnote[119]{Yi-pien 5408; 6385.} 119 it was also possible for a horizontal line of graphs to continue as a vertical column
or vice versa,\footnote[120]{Yi-pien 7456; Ping-pien 7.1–2.} 120
or even for an inscription to shift from vertical columns to a horizontal
column and back to a vertical column.\footnote[121]{Yi-pien 7767 (top of the right hyoplastron). On placement in general, cf. Serruys (1974), pp. 16-17, who diagrams the various possibilities. Li Ta-liang (1972), pp. 31-52, provides an exhaustive account of the various writing patterns on shell. For the symbols or \_, \ref{sec:chapters-ch05:5.6} (see see sec. 5.6).} 121 Some of these patterns tended to be associated
with specific plastron areas.\footnote[122]{Chou Hung-hsiang (1969), pp. 12-36, has made an exhaustive study of inscription patterns} 122
Yi-p'ing [1959], pp. 230, 233; but some graphic
inversions were clearly intentional [see \ref{ch:5} (see ch. 5), n.
34]). For crack numbers, wrongly written and
corrected, see n. 41. See too Chou Hung-hsiang
(1970), pp. 49-52, on the question of erasing
without reengraving and reengraving without
erasing for a particularly interesting case of
multiple erasures, see Kuo Mo-jo (1972), p. 3;
Ch’iu (1972), p. 43.
% source: scan 068, printed 51
Scapulas lacked the natural center line of plastrons. Graphs were generally written
in vertical columns; they were read downward, and the columns might move from the
center out toward the two edges or from the edges toward the center. 123 Generally, the
engravers cut the inscriptions into the region of the bone near the socket and along the
two sides, frequently placing them near the right edge of a left scapula and vice versa. 124
The direction of the cracks influenced the direction of the inscription. In general, the
inscriptions ran against the cracks. That is, if the pu crack faced right, then the columns
on various parts of the shell, both front and back.
He points out, for example, that in the region of
the epiplastron and entoplastron (fig. 3), more
pairs are inscribed running from the outside edge
toward the center e.g., Ping-pien 5.1-2; 7.1-2;
12.1-2; 14.1-2; Yi-pien 6385) rather than vice
versa (e.g., Ping-pien 276.3-4). On the basis of the
reconstituted Ping-pien plastrons, I would propose
that ideally (as in Ping-pien 1; 3; 5; 22; 34; 36;
37; 38; 43; 47, etc. multicolumn inscriptions on
the epiplastron were expected to run from the
outside toward the center; those on the entoplastron were to run from the center toward the
outside; on the hypoplastron and hyoplastron
they were expected to run from the center toward
the outside; and on the xiphiplastron, those near
the edge were expected to run toward the center,
while those near the spine were expected to run
toward the edge. Or, to put the rule less technically, on the top and bottom sections of the
plastron, those inscriptions placed near the edge
ran toward the center, those placed near the
center ran toward the edge; in the middle of the
plastron, inscriptions all ran toward the edge.
Fig. 17 illustrates this; for examples, see Ping-pien
7.1-2, 3-4; 12.1-2, 3-4; 14.1-2, 3-4, 5-6;
16.1-2, 3-4, 5-6, 7-8. But although this pattern
may have been the model, the diviner did not
always follow it; e.g., Chia-pien 221; 222 (see the
k'ao-shih); Ping-pien 10.1-2; 29.1-2; 30.2; 41.16–
17; 49.7; 276.1-2, 3-4, 5-6; cf. Ping-pien 611,
k'ao-shih, p. 95.
123. It was frequently the case that columns
on a left scapula ran to the right and vice versa;
e.g., Tung (1945), pt. 2, \ref{ch:5} (see ch. 5), pp. 20a-21b
(period V); ch. 6, pp. 6b, 7b (period II); ch. 10,
p. 2b (period II); Chia-t'u 110; 111; 124 (all period
III). But columns on left scapulas might also run
to the left and those on right scapulas to the right;
e.g., Tung (1945), pt. 2, ch. 6, p. 5a (period II);
ch. 7, p. 2a (period V); Chia-t'u 199 (period I).
And in some cases the columns on one scapula
might run in both directions; e.g., K'u fang 1595
(D) (period I), see Tung (1945), pt. 2, \ref{ch:3} (see ch. 3),
p. 28a; Chia-t'u 125 (period III); 198 (period I).
124. Chou Hung-hsiang (1969), pp. 37-52,
discusses the placement of inscription pairs on
scapulas; on pp. 52-57, he documents the similarity of pattern employed on both bone and
shell. The right and left placement of paired
charges requires more study. Preliminary research
suggests that in cases in which the diviner was
merely seeking information and had no preference
among the responses he might obtain, the charge
in the positive mode was generally placed on the
right side of the shell, the charge in the negative
mode on the left side, e.g., Ping-pien 28; 30. When
one response was preferred, however, the desired
alternative, whether couched in the positive or
negative mode, appears to have been placed on
the right, the undesired alternative on the left
(Shirakawa [1948], p. 36; Chang Ping-ch'üan
[1956], pp. 241–243; cf. Takashima [1973], p.
182; Serruys [1974], p. 25). Thus, the charge
7, ``Tzu Shang (?) will perhaps have
sickness'' (Chui-ho 292 [S449.1]), was expressed in
the positive mode; but since Tzu Shang was an
ally and perhaps a member of the royal family,
his possible sickness was undesirable, and the
inscription was carved on the left side of the shell.
The remains of the wished-for, negative charge
[], “Tzu [Shang] will not have sickness,''
can be discerned on the right side of the plastron.
The validity of this interpretation is supported by
the ``weakening particle'' ch'i *, which I translate
as ``perhaps'' (see \ref{ch:3} (see ch. 3), n. 41), on the undesir
able, left side. This hypothesis about right and
left placement, which is generally, though not
invariably, supported by the evidence (see, e.g.,
Ping-pien 1.3-4 [translated in\ref{sec:chapters-ch02:2.8} (see  sec. 2.8)]; 8 [fig. 8];
[fig. 12]; Yi-pien 3287 [S433.1]; 4130@@
[S245.4]), depends, of course, on the unverifiable
assumption that we can tell what the Shang king
was wishing for; for a case in which the hypothesis
fails, consider Ping-pien 235.1-2 (fig. 11), translated
in\ref{sec:chapters-ch02:2.8} (see  sec. 2.8)). For a convenient list of inscriptions
catalogued by the right-left position of positive-negative charges, see Chou Hung-hsiang (1969),
pp. 54-57; Li Ta-liang (1972), pp. 89-94. Note
that the practice of placing desired charges on the
right and undesired on the left of the front of the
bone or shell was apparently reversed on the back
of the shell (Chou Hung-hsiang [1969], pp. 55-
57), indicating that the bifurcation was thought
to inhere in the bone or shell itself, not in the
viewer.
% source: scan 069, printed 52
of the inscription moved from right to left, in the opposite direction; if the pu crack faced
left, then the inscription columns moved from left to right.\footnote[125]{Also note that if the pu crack itself faces right, then the pu graph associated with that crack frequently does so too, and vice versa (Ch'en Meng-chia [1956], p. 13); e.g., Ping-pien 10; 12.1-2, 5-6; 14.1-2; 16.1-2 (fig. 15); 18.1-2; 22.1-2, 5-6, 9-10\%; in this way, the mirror symmetry of the graph itself was preserved. But at other times symmetry was not preserved; e.g., Ping-pien 8.2 (fig. 8); 14.4-6; 16.6 (fig. 15); 20.2; 28.3. In these cases, all on the left side of the shell, the pu graph faces left, i.e., toward the edge, when one would expect it to face right, toward the spine, like the cracks themselves. This failure to observe symmetry, or to observe congruence between crack and graph, suggests that the left-facing pu on the right side of the shell was carved first, and that the engraver simply copied it on the left side of the shell without rotating the graph about its axis. This would accord with the view that the desired charge was usually carved on the right side (n. 124); one would expect the desired charge to have been carved first (cf. n. 130). Other graphs might also be reversed in this way, e.g., the graph for the diviner Ch'üeh in fig. 8 (cf. Serruys [1974], p. 18, who suggests ``that only a certain type and limited number of graphs are regularly written in mirror form'').} 125 Figure 8 is a typical example.
Exceptions to this rule, however, do exist.\footnote[126]{See, for example, Tung (1936), p. 115; (1945), pt. 2, ch. 6, pp. 5a, 7b; Chia-pien 221 and 222, k'ao-shih, pp. 35-36; Ping-pien 390; 611.} 126
In general, plastrons tended to be burned from top to bottom (sec. 1.6.3), with the
cracks and their inscriptions moving down the shell from the epiplastron or top of the
hyoplastron.\footnote[127]{E.g., Ping-pien 5; 8 (fig. 8).} 127 But the order was by no means assured. The sequence of inscriptions
might run from the epiplastron to the top of the hyoplastron, jump to the hypoplastron,
and then move back to the bottom of the hyoplastron;\footnote[128]{E.g., Ping-pien 22.} 128 and in some cases, the cracks or
inscriptions might move up the shell rather than down.\footnote[129]{E.g., Ping-pien 122. Li Ta-liang (1972), pp. 8-30, analyzes plastron use according to eleven different sequences.} 129 The inscriptions give some
indication that when a pair of charges was cracked, the right side of the shell was cracked
before the left.\footnote[130]{The inscription on the right tends to be complete, the charge on the left tends to be abbreviated (Chang Ping-ch'üan [1956], p. 244; cf. Chou Hung-hsiang [1969], pp. 60–62; Li Taliang [1972], p. 29); this suggests that the right side was given priority (cf. Serruys [1974], p. 29), a suggestion supported by the view that the favorable, desired charge was generally placed on the right; see nn. 124 and 125.} 130
In the case of scapulas, the right side of a left-hand scapula tended to be used first
and vice versa; and the burning of cracks, and thus the recording of associated charges,
tended to proceed either from bottom to top or, less frequently, from top to bottom.\footnote[131]{Tung (1945), pt. 2, \ref{ch:5} (see ch. 5), p. 22a; ch. 6, pp. 5a, 8a. Cf.\ref{sec:chapters-ch01:1.6.3} (see  sec. 1.6.3), par. 4.} 131
The bottom-to-top order was common in routine divinations such as those about misfortunes in the ten-day week, the night, or the royal hunt, in which the individual
charges ascended the bone as the days passed.\footnote[132]{E.g., Tung (1929), p. 75, figs. 37, 38; Shirakawa (1972), pp. 21, 30, figs. 3.1, 3.2, 3.8; cf. Matsumaru (1963), pp. 36, 51-52. For divinations running down the scapula, see Tung (1945), pt. 2, ch. 6, p. 7b; ch. 7, p. 2a.} 132 In another case, a series of divination
charges about the pin ritual zigzagged up the bone.\footnote[133]{E.g., Ts'ui-pien 176 (fig. 18).} 133 And in still other cases, particularly in periods I and IIa, the order was irregular.\footnote[134]{E.g., Chien-shou 29.5 (see Tung [1945], pt. 2, ch. 8, pp. 5a-b.} 134
It is clear from the above discussion that the engraver obeyed certain rules: the in-
% source: scan 070, printed 53
scriptions had to be near their cracks, and within one inscription, the graphs had to be
read from top to bottom or sideways; they were never read from bottom to top or in boustrophedon style. Further, the diviner's choice in patterns of placement and sequence was
limited. These considerations affect the modern reconstruction of scapulas and plastrons
from fragments; a putative matching piece and its inscriptions must fit one of the forms
described above (\ref{sec:chapters-ch05:5.7} (see see sec. 5.7)).
Engraving the Cracks
The pu cracks in scapulas, plastrons, and carapaces were themselves frequently
carved more deeply, but only in period I. 135 In some cases, the cracks were carved even
though no inscriptions were made; some cracks on a plastron might be carved, others
might not be; when the crack was carved, the inscription frequently straddled the crack
in a fan-chao situation. 136 This suggests, first, that the cracks were carved at the same time
as the inscription was engraved. 137 And it suggests, second, that the carving was done to
distinguish the cracks from the inscriptions by making them more visible. 138 This view is
supported by the fact that colored pigment (sec. 2.12) was sometimes added to these carved
cracks.
J
2. II
Engraving the Boundary Lines
The engravers also carved straight or sinuous boundary lines (modern term,
chieh-hsien) on the bone or shell (as in fig. 7). 139 These were presumably added at
the time the inscriptions were carved 140 and served to separate the divination sentences by
135. E.g., Ping-pien 57; Yi-pien 867; K'u-fang
(D). Tung (1949a), p. 261, cites other@@
examples on plastrons and carapaces. Ch'en
Meng-chia (1956), p. 16, notes that such carved
cracks were numerous on the shells excavated
from YH127; Tung (loc. cit.) says that some four
thousand fragments from this pit had such carved
cracks. He also proposes ([1945], pt. 2, \ref{ch:3} (see ch. 3),
p. 26b; [1954b], p. 356) that carved cracks are
characteristic of the reigns of Hsiao Yi and early
Wu Ting, but this is uncertain; we probably
have few, if any, inscriptions from the period
of Hsiao Yi (see \ref{ch:4} (see ch. 4), n. 24). Since, due to the
presence of the hollows, the shell was already
paper-thin where the cracks formed, it has been
suggested that a special knife, with a blade that
was semicircular in cross-section, may have been
used to cut the cracks so that the carver did not
break through the shell; this would explain the
semicircular cross-section and blunt ends of most
of these deepened cracks, though others still show
a V-form (Chou Hung-hsiang [1967/68], p. 27).
136. Tung (1945), pt. 2, \ref{ch:3} (see ch. 3), p. 26a; (1949a),
p. 261; he cites Ping-pien 96; 243; 502, etc. For
fan-chao, \ref{sec:chapters-ch02:2.4} (see see sec. 2.4).
137. This contrasts with the view of Chang
Ping-ch'üan (1967), p. 861, that the cracks were
probably engraved at about the same time as the
crack numbers and crack notations.
138. See \ref{ch:1} (see ch. 1), n. 93 on the tendency of cracks
to lose their visibility.
139. Tung (1929), p. 59, suggested that the
graph chao (small seal ) was a picture of
these boundary lines, but this is speculative.
140. The unique example of Ping-pien 390.5,
in which the boundary line, presumably as the
result of carelessness, separates the crack number
from its crack, suggests that the boundary lines
were carved after the crack numbers (cf. Chang
Ping-ch'üan [1956], p. 240). Li Ta-liang (1972),
p. 122, suggests that the divinatory *tieng was
omitted in two cases (Ping-pien 124.11; 306.9)
because the engraver encountered the boundary
line. If so, this would suggest that the boundary
% source: scan 071, printed 54
forming a series of compartments or ``spheres of influence'' when there was danger that
the graphs of one inscription might be confused with those of another. Some carved
boundaries apparently separated cracks alone,\footnote[141]{E.g., Ping-pien 519; 550; Yi-pien 1722; 3343. A pu crack on Yi-pien 9085 is surrounded by a boundary cartouche, but this may be only a doodle.} 141 but those which served as a boundary
between inscriptions were more common.\footnote[142]{E.g., Chia-pien 182; 690; 2868; 3313; Pingpien 126; 167; 243: 247 (fig. 12); 304; 323; 473; Yi-pien 145; 151; 811; 1988; Ching-hua 1; 2 (fig. 14); 3. For other examples, see Chang Ping-ch'üan (1954), p. 229; (1956), p. 240. Certain puzzling cases (see the k'ao-shih to Yi-chu 27 and 73 in which a boundary line seems to divide graphs in the same inscription may perhaps be explained as instances of two abbreviated inscriptions.} 142 The effort that must have been involved in
carving these lines suggests they served an important function.\footnote[143]{Cf. n. 102. Boundary lines might appear on several contemporary bones dealing with the same topic; cf., for example, H103: 18+20 (KK [1975.1], p. 42, plate 11.2), with Chih-hsü 64 and Ts'ui-pien 79 (both S466.3). On the diminishing use of boundary lines in the later periods, see \ref{ch:4} (see ch. 4), n. 88.} 143

\section[Pigmentation]{\origsecnum{2.12}Pigmentation}
\label{sec:chapters-ch02:2.12}

We have already noted (sec. 1.3.3) the care with which the oracle bones and
shells were smoothed and polished. The not infrequent addition of colored pigments,
placed in the grooves of the engraved graphs or in the pu cracks themselves, is further
evidence of the value placed on them.\footnote[144]{According to Britton (1937), p. 1, the colors also ``occur on the surface, as though originally spread over the entire bone or shell in the manner of a varnish.'' No other report refers to such surface colors. As Britton himself says on the same page, the colors in the cracks appear ``to have been applied with great care''; and he later wrote that the color in the grooves ``seems to have been applied in fluid form by brush, the brush point having been fine enough to line the grooves without smearing the margins'' (1940), p. 15. Conceivably, the pigments were rubbed into the grooves and the surface then wiped clean.} 144 This practice, like that of engraving the cracks
(sec. 2.10) was prevalent under Wu Ting; it appears to have been less common in later
periods.\footnote[145]{According to Chen Meng-chia (1956), p. 16, brownish-red pigment was still used in later periods, but having lost its color with age, it has frequently been overlooked. Britton (1940), p. 7, suggested that black pigment was used after the red phase of Wu Ting's reign; but black and red do appear on the same oracle bone; e.g., Ping-pien 208 (see the colored frontispiece to Ping-pien, vol. 2, pt. 2); Ning-hu 2.25 (D); 2.26 (D); 2.28-2.31 (D) (after Chen Meng-chia [1956], p. 15). According to Tung (1933), p. 421, the inscription of diviner Wei on Chia-pien 3339 was colored red, that of diviner Hsüan on the same bone was colored black (the k'ao-shih, p. 424, however, records only the red). For examples of pigmented bones and shells, see Chia-pien 50; 65; 67; 68; 84; 393; Ping-pien 12-21 (sec. 3.7); 90\%; 95\%; 97; 207; 209; 279; 281; 283; 284; Yi-pien 1877; 6274; Ching-hua 2 (fig. 14). Britton (1940), pp. 15-16, cites cases in the Princeton collection (Ch'i-chi, P). An asterisk precedes the rubbing numbers of pigmented bones and shells in Menzies, Ming-hou, and White. Li Yen (1970), plate 1, provides a not entirely satisfactory color reproduction of a scapula with red pigment in the Royal Ontario Museum.} 145
line had been carved prior to the inscription
itself; this would be particularly likely to have
been true in the case of Ping-pien 306.9, where the
adjacent inscription (306.17), i.e., the one that
would have taken up space, refers to a divination
made two days later that might not have been
present when 306.9 was carved. But more research
is needed; a fuller study of the boundary lines
depends on, and may help resolve, the questions
raised in n. 90.
% source: scan 072, printed 55
The Nature of the Colors
The colors have been described as brilliant red, vermillion, dark red tending to
black, brownish red, dark brown, and black.\footnote[146]{Britton (1937), p. 1; Chang Ping-ch'üan (1957), p. 4; (1967), p. 865. Li Yen (1966), p. 2, refers to a scapula smeared with yellow.}\footnote[147]{Chang Ping-ch'üan (1967), p. 865.}\footnote[148]{In some cases the colors have faded badly; e.g., Ping-pien 59; 69; 78; 90; k’ao-shih, pp. 94, 103, 111, 122. The distribution of brown and black found on Ping-pien 12-21 (see table 8. makes it likely that the two colors, in this case at least, were originally the same.}\footnote[149]{Tung (1948a), p. 127, referring to Chiapien 3056; cf. Tung (1965), p. 47, on Ching-hua 2 (fig. 14). He made a similar comment about Ping-pien 284: ``the graphs all adorned with cinnabar, red, as fresh as new'' ([1954b], p. 354), but according to the k'ao-shih, p. 348, the graphs were filled with black.}\footnote[150]{Benedetti-Pichler (1937), p. 152. The red pigment was taken from Ch'i-chi, P 39 (D), the black from P 13 (D) (ibid., p. 150). Britton, apparently extrapolating from Chou ritual practice, told Benedetti-Pichler that ``the black pigment is thought to be the blood of animals used in divination'' (ibid., p. 149. As Benedetti-Pichler notes (p. 152), the negative hemin test does not conclusively exclude the presence of blood. As to the red coloring, Rutherford J. Gettens subsequently concluded that it ``was made from natural cinnabar, because of the heterogeneous size of the particles'' (Britton [1940]. p. 7).}\footnote[151]{Lin Sheng (1964 pp. 99, 101. Ch'en Meng-chia (1956), p. 16, identifies the black as carbon, but I do not know on what evidence.}\footnote[152]{Chang Ping-ch'üan 1967), p. 865.}\footnote[153]{According to Chien Meng-chia (1956), p. 15, citing five examples, large characters were filled with red pigment, small ones on the same bone or shell with black (cf. Chu Fang-p'u [1935], ch. 6, p. 18a; Kaizuka [1967], p. 204). Chang Ping-ch'üan (1967), p. 866, however, finds no relation between the color and size of the graphs. Since some of the inscriptions which were originally colored are thought to have lost all their pigmentation, it is unlikely that we shall be able to reconstruct the Shang situation. CEO E ]}
% source: scan 073, printed 56
Quite possibly, there was no fixed rule, the diviner using whichever color was available
and only when it was available. A majority of the inscriptions, in any event, are not thought
to have been colored.\footnote[154]{Given the tendency of pigment to fade, and of the medium (such as gum or resin?) with which it may originally have been mixed to decay (Goodyear [1971], pp. 77, 95, 126-127), the possibility that all graphs were once colored cannot be entirely excluded. It is generally thought, however, that this was not the case, even in period I (e.g., Tung [1949a], p. 265; Ping-pien 17.1-2, k'ao-shih, p. 40). ``} 154
By analogy with the Yünnanese case, in which the modern diviner rubbed ash into
the cracks, apparently to make them more visible, it is probable that the Shang used
pigments for the same reason.\footnote[155]{See the sources cited in n. 151; see too \ref{ch:1} (see ch. 1), n. 98. Britton (1940), p. 7: ``Glyphs cut in fresh bone are scarcely visible under well-diffused light and this probably explains the use of colour in the glyphs.'' Obscure references to inking the shell appear in Chou-li (``Ch'un-kuan, Pu-shih''; CLCY, ch. 47, p. 14b; Biot [1969], vol. 2, p. 75; cf. Eisenberger [1938], p. 67), Li-chi (``Yü-tsao,' p. 166; Couvreur [1950], vol. 1, pp. 682-683), and in the pseudo-K'ung commentary to Shangshu, ``Lo-kao'' (Shih-san-ching chu-shu, ch. 15, p. 16a; Legge [1865], p. 437, n.). But these passages, which have usually been taken to refer to the inking of pseudocracks on the shell and to the ability of the fire to dry the ink or to produce cracks replicating the inked ones, cannot be used to understand the Shang evidence; rather, the opposite is true (Lin Sheng [1964], pp. 100-101; cf. Chavannes [1911], p. 132; Jao [1961], p. 953). That the primary motive may have been visibility is suggested by the following comment on classical Greek inscriptions: ``For greater clarity, stone-cut letters were often also coloured with paint, as is in fact not infrequent in modern times. The colours have seldom survived, but where traces do remain they suggest that red was the colour ordinarily used for the purpose. Where two colours were used, alternate lines were painted in red and black. This two-colour system may have been more widespread than the scanty remains of it so far found suggest'' (Woodhead [1967], p. 27; cf. p. 129, n. 1).} 155 We can only speculate that the coloring may also have
served to beautify the graphs, or to sanctify and render them more potent.\footnote[156]{The hematite surrounding the skeletons in the upper cave at Chou-k`ou-tien (Chang Kwang-chih [1968], p. 70) and at Yang-shao sites such as Ta-ho-ts'un and the fact that the bones of some of the persons buried at Hsiao-t'un were colored red, a frequent occurrence in Shang graves (Shih Chang-ju [1952], pp. 466, 468–485), indicate that red, the color of blood, may have had religious significance in both Neolithic and Shang times (cf. Kaizuka [1967], pp. 204-205). Shirakawa (1972), p. 29, points out that the ming-ch'i in the Shang royal tombs were frequently daubed with red, probably to sanctify them. I am dubious about the potency argument because \listitem{1} we find that the marginal notation ``Ch'üeh sent in 150,'' carved on the bridge of a plastron, was also smeared in red (e.g., Ping-pien 95.2; Yi-pien 6423 [after Tung (1948a), p. 124]); so were boundary lines (Chia-pien 65, k'ao-shih, p. 14), and it is hard to see what function could have been served by rendering the notation of provenance, or a boundary line, more potent; \listitem{2} in complementary charge pairs, such as ``Sick tooth will be favorable will perhaps not be favorable'' (Pingpien 16.9-10 [fig. 15]; \ref{sec:chapters-ch03:3.7.1.2} (see see sec. 3.7.1.2)), it is hard to see why the Shang would have wanted to make the undesirable alternative as potent as the desirable. The potency argument would be strengthened if it could be shown that there was a tendency for prognostications to be colored in this way, i.e., as a sign that the king was attempting to make his forecasts come true. But this does not appear to have been the case.} 156 Whatever
the function, a blend perhaps of the practical, aesthetic, and religious, it appears to have
been an embellishment, not an essential feature, of Shang divination practice.
